\documentclass[11pt]{ctexart}
\usepackage[a4paper,margin=1in]{geometry}
\usepackage{graphicx}
\usepackage{xcolor}
\usepackage{hyperref}
\definecolor{refgray}{gray}{0.45}
\title{\textbf{Market Views｜全球投行叙事汇编}\\[0.4em]{\large AI变现进入验证期，消费与工业分化加剧，存储定价强势延续}}
\author{KC桌面 自动生成}
\date{260731}
\begin{document}
\maketitle
\tableofcontents
\newpage
\section{一页摘要}
\begin{itemize}
  \item 全球权益资金流入创历史新高，但零售端情绪已明显降温，增量资金主要来自机构被动配置。
  \item AI货币化仍处早期，但微软Azure增速连续加速至43\%，M365 Copilot付费席位单季净增1000万，验证企业级AI订阅需求。
  \item 存储（DRAM/HDD）定价在2026年下半年保持双位数季度环比增长，供应短缺是核心驱动力，UBS预计DRAM位元需求增速从2026年22\%升至2027年36\%。
  \item 中国消费市场结构性分化加剧，高净值客群消费韧性显著强于大众消费者，奢侈品、高端运动品牌均呈现这一趋势。
  \item 全球电网投资持续高景气，变压器等关键设备供需紧张，中国供应商凭借供应链优势填补缺口。
  \item 美国市场定价环境稳定，为高敞口日系车企提供盈利支撑；中国乘用车市场新车发布效应消退，订单分化加剧。
  \item 跨国药企在华业绩因产品组合与政策暴露度不同而显著分化，创新产品成为增长关键。
  \item 科技与半导体板块的去杠杆已基本完成，进一步下行空间有限，逆周期因子仍在中间价形成中发挥缓冲作用。
\end{itemize}
\section{全球投行叙事汇编}
今日纳入 63 篇投行/券商研究，按 12 个市场、资产与产业主题系统整理。
\newpage
\subsection{宏观与策略：分化与再定价}
\textbf{全球市场正经历资金结构、消费行为与产业周期的三重分化：权益资金流入创纪录但零售情绪降温，美国消费者转向价值导向，中国消费呈现K型复苏，而存储与数据中心等AI相关领域则迎来结构性需求上修。}
\subsubsection*{主流共识}
\begin{itemize}
  \item 全球权益资金流入创历史新高，但零售端情绪已明显降温，增量资金主要来自机构被动配置。
  \item 美国消费者支出整体保持韧性，但系统性向价值导向迁移，低收入群体承压最大。
  \item 中国消费市场呈现K型分化，高端消费韧性明显，大众可选消费恢复仍面临阻力。
\end{itemize}
\subsubsection*{分歧与条件差异}
\begin{itemize}
  \item 对AI资本开支周期的持续性存在分歧：Bernstein认为Equinix上调资本开支至50-70亿美元/年（开发回报约25\%）是结构性机会，而市场担忧资本效率与现金流滞后。
  \item 存储周期见顶与否：UBS预计DRAM位元需求增速从2026年22\%升至2027年36\%，与市场对周期见顶的直觉相反。
  \item 中国生猪养殖周期拐点时间：MS将周期性复苏推迟至2027年三季度，因生产效率提升延缓了供给收紧。
\end{itemize}
\subsubsection*{机构观点}
\paragraph{BARC} 7月全球权益净流入1290亿美元，年初至今累计6590亿美元创纪录，但零售情绪已明显降温（AAII牛熊指标降至负一个标准差，恐惧贪婪指数回落）。对冲产品和CTA多头重新累积，整体权益持仓接近历史高位，共同基金现金水平处于十年底部，市场应对冲击的缓冲空间有限。欧洲因对科技/AI敞口较低、价值因子持续跑赢，在不确定性上升中相对有利。
{\small\color{refgray}数据：7月全球权益净流入1290亿美元\par}
{\small\color{refgray}数据：年初至今累计流入6590亿美元，创纪录\par}
{\small\color{refgray}数据：AAII牛熊指标降至负一个标准差\par}
{\small\color{refgray}数据：共同基金现金水平处于十年底部\par}
{\small\color{refgray}边际变化：此前市场关注资金流入总量，现转向关注资金结构变化：零售端降温、对冲基金重新加仓、持仓接近历史高位。\par}
\paragraph{BARC} 梅赛德斯-奔驰Q2工业业务净流动性304亿欧元，但自由现金流同比下滑41\%至11亿欧元，管理层预计下半年乘用车利润率弱于上半年。中国业务利润贡献下降（8.44亿欧元估值调整，其中7.52亿为联营公司减值），叠加原材料成本上升与激进股东回报，信用质量在评级区间内偏弱。
{\small\color{refgray}数据：工业业务净流动性304亿欧元\par}
{\small\color{refgray}数据：自由现金流同比下滑41\%至11亿欧元\par}
{\small\color{refgray}数据：中国业务估值调整8.44亿欧元\par}
{\small\color{refgray}数据：H1股东回报50亿欧元\par}
{\small\color{refgray}边际变化：市场此前定价未充分反映中国业务利润贡献下降、原材料成本上升与激进股东回报三者叠加的影响。\par}
\paragraph{Bernstein} 中国Q2 GDP增速从5.0\%回落至4.3\%，零售销售仅增长0.2\%，消费呈现K型分化。高端白酒韧性突出（茅台飞天批价较底部涨13\%），运动服饰线上增长10\%但线下疲弱。酒店高端与经济型RevPAR强劲，中间段承压。消费者信心虽改善但进程缓慢，大件商品和汽车销售仍在调整。
{\small\color{refgray}数据：Q2 GDP增速4.3\%\par}
{\small\color{refgray}数据：零售销售增长0.2\%\par}
{\small\color{refgray}数据：茅台飞天批价较底部涨13\%\par}
{\small\color{refgray}数据：运动服饰线上销售增长10\%\par}
{\small\color{refgray}边际变化：此前市场关注整体消费复苏，现转向K型分化：高端消费韧性明显，大众可选消费恢复阻力仍大。\par}
\paragraph{Bernstein} Equinix Q2业绩超预期，将年度AFFO每股增长指引从5-9\%上调至9-12\%，资本开支从30-40亿美元提升至50-70亿美元。约80\%投入25个核心市场，开发回报约25\%，无需股权融资。AI需求已嵌入核心业务：前十大AI模型提供商中八家使用其网络，新增9700个互联节点创纪录。
{\small\color{refgray}数据：AFFO每股增长指引从5-9\%上调至9-12\%\par}
{\small\color{refgray}数据：资本开支从30-40亿提升至50-70亿美元/年\par}
{\small\color{refgray}数据：开发回报约25\%\par}
{\small\color{refgray}数据：新增9700个互联节点创纪录\par}
{\small\color{refgray}边际变化：此前市场关注数据中心资本开支激进风险，现转向AI需求结构性嵌入与资本效率优势。\par}
\paragraph{MS} 中国生猪养殖产能出清速度慢于预期，生产效率提升（MSY提高）抵消了能繁母猪存栏下降对供给的影响。预计2026年行业生猪均价11元/公斤（同比-23\%），周期性复苏推迟至2027年三季度。牧原股份完全成本11.5元/公斤，MSY 25 vs 行业平均20，成本优势在周期底部依然显著。
{\small\color{refgray}数据：2026年行业生猪均价预计11元/公斤\par}
{\small\color{refgray}数据：同比下跌23\%\par}
{\small\color{refgray}数据：牧原完全成本11.5元/公斤\par}
{\small\color{refgray}数据：MSY 25 vs 行业平均20\par}
{\small\color{refgray}边际变化：此前市场预期2026年下半年迎来周期拐点，现推迟至2027年三季度，因生产效率提升延缓供给收紧。\par}
\paragraph{UBS} SK海力士自6月高点回调，但基本面未变。DRAM位元需求增速预计从2026年22\%升至2027年36\%，NAND从20\%升至23\%。长期协议（LTA）已签10份，短期ASP承压但长期盈利稳定性改善。2027-2028年EPS下调19\%但仍高于市场共识17\%以上。预计2026-2028年自由现金流分别达188/320/374万亿韩元。
{\small\color{refgray}数据：DRAM位元需求增速2027年36\%\par}
{\small\color{refgray}数据：NAND增速2027年23\%\par}
{\small\color{refgray}数据：LTA已签10份\par}
{\small\color{refgray}数据：2026-2028年FCF 188/320/374万亿韩元\par}
{\small\color{refgray}边际变化：此前市场关注DRAM短期ASP与资本开支上升，现转向智能体AI对存储需求的结构性拉动与LTA改善盈利稳定性。\par}
\paragraph{Bernstein} 奔驰Q2调整后EBIT 22.99亿欧元（同比+16\%），超卖方共识40\%，但乘用车报告净利润仅4900万欧元（同比-94\%），因中国合资公司减值7.96亿欧元。排除中国后全球销量同比+2\%，BEV全球销量+51\%（欧洲+87\%）。管理层计划将低成本地区产能弹性提升至40万台，成本较德国低70\%。
{\small\color{refgray}数据：调整后EBIT 22.99亿欧元，超预期40\%\par}
{\small\color{refgray}数据：乘用车报告净利润4900万欧元，同比-94\%\par}
{\small\color{refgray}数据：中国合资公司减值7.96亿欧元\par}
{\small\color{refgray}数据：BEV全球销量+51\%\par}
{\small\color{refgray}边际变化：此前市场关注调整后利润韧性，现转向中国业务减值对报告利润的实质性影响与成本削减加速。\par}
\paragraph{Bernstein} 美国消费者支出整体保持韧性，但系统性向价值导向迁移。沃尔玛、好市多、美元店受益，家居改善板块承压。餐饮业面临成本与需求双重挤压，高端与平价两端分化。功能饮料赛道表现突出，Celsius和Keurig Dr Pepper组合强劲。通胀不确定性将持续，消费者价格敏感度不会很快消退。
{\small\color{refgray}数据：通胀率维持在3.5\%以上\par}
{\small\color{refgray}数据：沃尔玛、好市多、美元店客流增长\par}
{\small\color{refgray}数据：牛肉成本上涨\par}
{\small\color{refgray}数据：功能饮料赛道增长突出\par}
{\small\color{refgray}边际变化：此前市场关注消费支出总量韧性，现转向结构性变化：消费者系统性地向价值导向迁移，不同收入层级和品类分化加剧。\par}
\subsubsection*{关键数据}
\begin{itemize}
  \item 7月全球权益净流入1290亿美元，年初至今累计6590亿美元创纪录
  \item 奔驰中国合资公司减值7.96亿欧元，乘用车报告净利润同比-94\%
  \item Equinix资本开支从30-40亿提升至50-70亿美元/年，开发回报约25\%
  \item DRAM位元需求增速预计从2026年22\%升至2027年36\%
  \item 中国生猪养殖周期拐点推迟至2027年三季度，2026年均价预计11元/公斤
  \item 美国通胀率维持在3.5\%以上，消费者系统性转向价值导向
  \item 茅台飞天批价较底部涨13\%，运动服饰线上销售增长10\%
\end{itemize}
\begin{figure}[htbp]
\centering
\includegraphics[width=0.88\linewidth]{figures/fig_001.jpg}
\caption{FIGURE 1 - \#\# Relative value - Underweight From a credit perspective, the balance sheet remains strong, with €30.4bn of net industrial liquidity after €5bn of shareholder distributions in H1. However, management expects H2 Cars profitability to be weaker than H1 due to h}
\end{figure}
\begin{figure}[htbp]
\centering
\includegraphics[width=0.88\linewidth]{figures/fig_007.jpg}
\caption{FIGURE 1 - \#\# Equity positioning near the highs \#\# Long only inflows were elevated again in July, albeit skewed by record China inflows FIGURE 2. Ytd equity inflows are now tracking ahead of record high bar set in 2021, at \textbackslash{}\$659bn FIGURE 3. Asset manager net futures posi}
\end{figure}
\begin{figure}[htbp]
\centering
\includegraphics[width=0.88\linewidth]{figures/fig_012.jpg}
\caption{EXHIBIT 2 - EXHIBIT 2: China GDP grew by 4.5\% YoY in Q2 EXHIBIT 3: The GDP deflator improved from -0.6\% in Q4 to -0.1\% in Q1 EXHIBIT 4: Export growth picked up in Q2}
\end{figure}
\begin{figure}[htbp]
\centering
\includegraphics[width=0.88\linewidth]{figures/fig_017.jpg}
\caption{EXHIBIT 3 - EXHIBIT 3: EQIX Long Term Guidance EQIX Revenue Growth by Segment EXHIBIT 4: EQIX Revenue YoY Growth by Segment (\%) EXHIBIT 5: EQIX EBIT Margin EXHIBIT 6: EQIX Diluted EPS (\textbackslash{}\$, \%)}
\end{figure}
{\small\color{refgray}本节综合 8 篇研究；涉及 BARC、Bernstein、MS、UBS。}
\newpage
\subsection{AI生态演进：从模型竞赛到基础设施与货币化验证}
\textbf{AI产业正从模型参数竞赛转向生态构建与货币化验证阶段，企业采购重心从‘选模型’转向‘选生态’，边缘推理、先进封装与SaaS AI增量成为关键验证点。}
\subsubsection*{主流共识}
\begin{itemize}
  \item AI货币化仍处早期，尚未成为收入加速的驱动力，SaaS板块估值折价45\%但缺乏重估催化剂。
  \item 先进封装设备订单持续回暖，TCB与光通讯工具受数据中心CPO与HBM需求驱动。
  \item 企业AI采购从模型选择转向生态构建，微软Azure加速增长验证平台粘性。
\end{itemize}
\subsubsection*{分歧与条件差异}
\begin{itemize}
  \item AI对SaaS收入贡献是‘增量’还是‘拐点’：Figma等公司AI信用货币化初现，但多数公司贡献仍为低个位数。
  \item 中国大模型竞争格局：智谱ARR提前达标但面临Kimi K3参数跃升压力，规模定律是否持续有效存分歧。
  \item 边缘AI部署节奏：Cloudflare渠道调研显示份额领先，但大规模部署预计在2027年，短期收入贡献有限。
\end{itemize}
\subsubsection*{机构观点}
\paragraph{GS} ASMPT订单连续两季环比增长45\%和24\%，3Q26预期高个位数环比增长，驱动力来自TCB与光通讯工具，受益于数据中心CPO采用率上升、高端CPU/HBM需求及先进封装向面板级迁移。2026年营收预测上调至194亿港元，毛利率从40.3\%上调至41.6\%，但净利润预测基本不变因税率与一次性影响。
{\small\color{refgray}数据：1Q26订单环比+45\%，2Q26环比+24\%\par}
{\small\color{refgray}数据：3Q26订单预期高个位数环比增长\par}
{\small\color{refgray}数据：2026年营收预测从181.51亿上调至194亿港元\par}
{\small\color{refgray}数据：毛利率从40.3\%上调至41.6\%\par}
{\small\color{refgray}边际变化：订单增长驱动力从传统封装转向TCB与光通讯，反映AI封装需求结构性变化。\par}
\paragraph{GS} 微软4QFY26营收900亿美元（同比+17\%），EPS 4.74美元（同比+23\%），Azure增速43\%超预期，产能受限下仍加速增长。Copilot付费用户突破3000万，单季净增1000万。企业AI采购重心从‘选模型’转向‘选生态’，微软角色更趋战略性，GPU上线时间缩短近50\%，AI单位经济改善。
{\small\color{refgray}数据：4QFY26营收900亿美元，同比+17\%\par}
{\small\color{refgray}数据：Azure增速43\%（固定汇率），超市场预期40\%\par}
{\small\color{refgray}数据：Copilot付费用户突破3000万，单季净增1000万\par}
{\small\color{refgray}数据：GPU上线时间缩短近50\%\par}
{\small\color{refgray}边际变化：微软从模型提供商转向生态中枢，AI货币化路径从单一模型转向多模型自动匹配与定价改进。\par}
\paragraph{HSBC} 智谱AI 2026年7月ARR达10亿美元，比指引提前5个月，HSBC上调年底ARR至20亿美元，2027年底至65亿美元。但收入结构向低毛利API倾斜，竞争加剧致定价承压，H股增发后股本扩张。Kimi K3参数从1万亿跃升至2.8万亿，验证规模定律，智谱GLM-5.2仅7440亿参数，追赶窗口收窄。
{\small\color{refgray}数据：2026年7月ARR达10亿美元，提前5个月\par}
{\small\color{refgray}数据：2026年底ARR预测从10亿上调至20亿美元\par}
{\small\color{refgray}数据：Kimi K3总参数2.8万亿，激活参数1040亿\par}
{\small\color{refgray}数据：智谱GLM-5.2总参数7440亿\par}
{\small\color{refgray}边际变化：中国大模型竞争从参数竞赛转向ARR验证，但低毛利API占比上升稀释利润率。\par}
\paragraph{MS} 金山办公2Q26收入预计同比增长18-33\%，对应16-18亿元，增长来自WPS个人订阅与WPS 365企业协作。调整后利润增长3-45\%，中位数约24\%与收入增速持平，利润率稳定。市场对AI冲击传统办公软件的担忧缓解，但估值偏高（2026年PE 63倍）是主要风险。
{\small\color{refgray}数据：2Q26收入预计16-18亿元，同比+18-33\%\par}
{\small\color{refgray}数据：1H26收入预计同比+21-28\%\par}
{\small\color{refgray}数据：调整后利润中位数同比+24\%\par}
{\small\color{refgray}数据：2026年PE 63倍\par}
{\small\color{refgray}边际变化：AI担忧缓解，收入增长来自存量用户付费转化与新场景拓展，利润率稳定性成关键信号。\par}
\paragraph{MS} SaaS应用软件2Q26需求基本面稳固，但AI货币化仍处早期，管理层维持谨慎展望，广泛上调修正概率低。板块EV/S倍数低于五年均值，相对软件板块折价45\%。Figma AI信用货币化首个完整季度，超75\%组织用户继续消费；HubSpot AI采用率约20-25\%，信用额度贡献低个位数。
{\small\color{refgray}数据：SaaS应用EV/S倍数低于五年均值\par}
{\small\color{refgray}数据：相对软件板块折价45\%\par}
{\small\color{refgray}数据：Figma超75\%组织用户继续消费AI信用额度\par}
{\small\color{refgray}数据：HubSpot AI采用率约20-25\%\par}
{\small\color{refgray}边际变化：AI货币化从概念验证进入早期收入贡献，但尚未成为收入加速驱动力，板块重估需等待明确催化剂。\par}
\paragraph{MS} Cloudflare正从CDN/安全公司演变为边缘AI推理基础设施关键参与者。渠道调研显示其是过去三个月边缘推理市场份额增长最快厂商，55\%合作伙伴认为其份额增长最快，63\%选为边缘AI计算首选平台。Q2营收预计约6.65亿美元（同比+29.8\%），cRPO同比增长预计达33\%。
{\small\color{refgray}数据：55\%合作伙伴认为Cloudflare边缘推理份额增长最快\par}
{\small\color{refgray}数据：63\%选为边缘AI计算首选平台\par}
{\small\color{refgray}数据：Q2营收预计约6.65亿美元，同比+29.8\%\par}
{\small\color{refgray}数据：cRPO同比增长预计达33\%\par}
{\small\color{refgray}边际变化：边缘AI从概念走向早期收入贡献，Cloudflare在分布式推理领域建立领先地位，但大规模部署预计在2027年。\par}
\paragraph{UBS} ASEAN市场技术板块仅占MSCI ASEAN权重4\%，却贡献指数今年至今25\%回报。55\%成份股估值仍低于10年均值，指数被少数权重股拉高。外资流向分化：泰国和马来西亚获净买入，印尼和菲律宾持续流出。AI直接敞口有限，但数据中心、电力等供应链环节仍有机会。
{\small\color{refgray}数据：技术板块占MSCI ASEAN权重4\%\par}
{\small\color{refgray}数据：贡献指数今年至今25\%回报\par}
{\small\color{refgray}数据：55\%成份股估值低于10年均值\par}
{\small\color{refgray}数据：泰国和马来西亚获外资净买入\par}
{\small\color{refgray}边际变化：AI主题持仓集中度提升，分散化需求可能成为ASEAN市场重新获得关注的催化剂。\par}
\paragraph{Bernstein} Air Liquide 在2026年上半年交出了一份让市场感到复杂的成绩单。Bernstein分析师在最新报告中指出，这份业绩更像是一个“减速带”——短期数据令人困扰，但公司的长期基本面并未受到实质性影响。报告的核心判断是：1H26的业绩波动为读者提供了一个观察窗口，而非趋势转折点。
{\small\color{refgray}数据：这份研报的独特价值在于，它没有停留在对短期业绩的简单解读，而是深入拆解了Air Liquide增长引擎的运转状态。报告认为，公司的订单积压仍处于历史高位，电子和工业商业务有望在下半年加速增长，而10月5日的资本市场日将成为重要的催化剂事件。\par}
{\small\color{refgray}数据：报告中最具说服力的证据来自订单积压数据。截至2026年第二季度，Air Liquide的订单积压达到创纪录的60亿欧元，环比新增5亿欧元。这一数字意味着公司未来2-3年的收入可见度极高，是判断增长可持续性的核心指标。\par}
\subsubsection*{关键数据}
\begin{itemize}
  \item ASMPT 1Q26订单环比+45\%，2Q26环比+24\%
  \item 微软Azure增速43\%超预期，Copilot付费用户突破3000万
  \item 智谱AI ARR提前5个月达10亿美元，Kimi K3参数跃升至2.8万亿
  \item SaaS应用EV/S折价45\%，AI货币化贡献仍为低个位数
  \item Cloudflare 55\%合作伙伴认为其边缘推理份额增长最快
  \item ASEAN技术板块4\%权重贡献25\%回报，55\%成份股估值低于均值
\end{itemize}
\begin{figure}[htbp]
\centering
\includegraphics[width=0.88\linewidth]{figures/fig_079.jpg}
\caption{Exhibit 1 - Exhibit 1: ASMPT's BB ratio was 1.43 by 2Q26 (vs. 1.43 by 1Q26) Exhibit 2: ASMPT's revenues: Recovering trend Earnings revision: We factor in ASMPT's 2Q26 result and raise our 2026E revenue and GM to reflect the stronger the ex}
\end{figure}
\begin{figure}[htbp]
\centering
\includegraphics[width=0.88\linewidth]{figures/fig_111.jpg}
\caption{Exhibit 1 - Exhibit 1: Comparison between Z.ai and Moonshot on key metrics}
\end{figure}
\begin{figure}[htbp]
\centering
\includegraphics[width=0.88\linewidth]{figures/fig_126.jpg}
\caption{Exhibit 1 - Exhibit 1: Overall Software Multiples and SaaS Applications Group Both Trading Below 5-Year Trailing Average Exhibit 2: SaaS Applications Software Group EV/S Multiples Trading at a 45\% Discount to Broader Software Exhibit 3: EV}
\end{figure}
\begin{figure}[htbp]
\centering
\includegraphics[width=0.88\linewidth]{figures/fig_133.jpg}
\caption{Exhibit 5 - Exhibit 5: Our June Survey Shows 84\% of Resellers In line or Above Target, Slightly Below Last Quarter}
\end{figure}
{\small\color{refgray}本节综合 8 篇研究；涉及 Bernstein、UBS、GS、HSBC、MS。}
\newpage
\subsection{半导体：AI驱动需求分化，存储与被动元件定价强势}
\textbf{AI基础设施投资正从GPU向存储、测试设备和被动元件扩散，但不同细分市场的定价韧性和需求可见性出现显著分化：HDD和DRAM因供应受限而定价持续走强，而模拟芯片和部分设备商则面临产品组合调整的压力。}
\subsubsection*{主流共识}
\begin{itemize}
  \item AI推理芯片的测试需求正在成为设备市场新的增长引擎，而非仅仅依赖GPU测试时间的延长。
  \item 存储（DRAM/HDD）定价在2026年下半年保持双位数季度环比增长，供应短缺是核心驱动力。
  \item 被动元件（MLCC/钽电容）产能利用率逼近满载，AI相关营收占比持续提升。
  \item 半导体设备市场2026年WFE预测上调至1500亿美元以上，先进封装增速显著高于整体。
\end{itemize}
\subsubsection*{分歧与条件差异}
\begin{itemize}
  \item HDD需求可见性已延伸至2028年，而模拟芯片（如NXP）的汽车终端市场出现调整信号，数据中心业务占比仅3\%。
  \item KLA的2026年业绩受产品组合影响成为份额让渡方，但MS认为2027年工艺控制强度回升将推动其跑赢WFE。
  \item HBM定价上行空间存在分歧：GS预计三星电子HBM定价明年同比增长87\%，高于彭博一致预期的52\%。
  \item 意法半导体三季度指引低于预期，但订单出货比接近2且分销库存低于目标，复苏节奏而非趋势中断。
\end{itemize}
\subsubsection*{机构观点}
\paragraph{Bernstein} NXP Q2业绩扎实，营收34.96亿美元、EPS 3.61美元均小幅超预期，通信基础设施和工业物联网是主要超预期来源。但汽车终端市场尤其是中国出现调整信号（乘用车零售同比-20\%），数据中心业务占比仅约3\%，远低于多数同行。估值在板块回调中变得有吸引力，但上行空间取决于市场如何定价这些结构性矛盾。
{\small\color{refgray}数据：Q2营收34.96亿美元（预期34.63亿），EPS 3.61（预期3.54）\par}
{\small\color{refgray}数据：渠道库存稳定在11周，处于长期目标水平\par}
{\small\color{refgray}数据：数据中心业务约5亿美元，占收入3\%\par}
{\small\color{refgray}数据：中国乘用车零售量同比-20\%\par}
{\small\color{refgray}边际变化：Q3指引小幅高于预期，但汽车业务平稳与中国市场疲软形成对比，分析师对汽车和AI故事持保留态度。\par}
\paragraph{MS} Seagate F4Q26财报显示HDD定价是驱动收入、利润率、EPS和FCF超预期的核心因素，且定价上行空间尚未见顶。需求可见性从时间（近线产能分配延伸至2028年）和客户（从传统云扩展至AI模型开发者、机器人等）两个维度同时扩展。增量毛利率从72\%升至83\%，规模效应显现。
{\small\color{refgray}数据：9月季度定价同比+25\%，6月季度为+10\%\par}
{\small\color{refgray}数据：毛利率从47\%提升至52.7\%，下季度指引约57.5\%\par}
{\small\color{refgray}数据：CY27 EPS从市场共识\$35.75上调至\$59.33\par}
{\small\color{refgray}数据：增量毛利率83\%（上季度72\%）\par}
{\small\color{refgray}边际变化：需求可见性从“贯穿2027年”延伸到“2028年”，客户类型从传统云向新云、模型开发者、机器人等领域扩展，HDD需求基础正在拓宽。\par}
\paragraph{MS} KLA 2026财年业绩呈现混合信号，9月季度指引略低于预期，产品组合对工艺控制强度构成影响，KLA成为份额让渡方。但2027年结构性故事更清晰：125亿美元RPO（FY25底78.6亿）支撑，1D/1.4nm节点和先进封装（混合键合、SoIC）将推动工艺控制强度回升，预计2027年营收增长36\%，跑赢WFE的31\%。
{\small\color{refgray}数据：2026年WFE预测从1400亿美元以上上调至1500亿美元以上\par}
{\small\color{refgray}数据：2026年先进封装收入预期从10亿美元上调至11亿美元，同比+70\%\par}
{\small\color{refgray}数据：RPO达125亿美元（FY25底78.6亿）\par}
{\small\color{refgray}数据：预计2027年营收增长36\%，跑赢WFE的31\%\par}
{\small\color{refgray}边际变化：短期业绩受产品组合影响平淡，但2027年工艺控制强度回升的底层逻辑更清晰，RPO大幅增长提供可见度。\par}
\paragraph{GS} 爱德万测试Q1营业利润1900亿日元，超GS预期25\%，全年指引从6275亿日元大幅上调至8460亿日元。AI推理芯片的测试需求成为新增长引擎，SoC测试设备市占率从66\%持续上升，中期目标至少70\%。产能扩张加速，SoC测试机产能目标从5000台扩至10000台且时间点提前。
{\small\color{refgray}数据：Q1营业利润1900亿日元，超GS预期25\%\par}
{\small\color{refgray}数据：全年营业利润指引从6275亿上调至8460亿日元\par}
{\small\color{refgray}数据：SoC测试设备市占率从66\%持续上升，中期目标至少70\%\par}
{\small\color{refgray}数据：SoC测试机产能目标从5000台扩至10000台，时间点提前\par}
{\small\color{refgray}边际变化：AI推理芯片测试需求成为新增长引擎，而非仅依赖GPU测试时间延长；产能扩张加速，客户需求可见性约18个月。\par}
\paragraph{GS} MARUWA Q1营业利润63亿日元低于预期，但全年指引从297亿上调至337亿日元，核心驱动力是AI应用订单确认和交付日期锁定。Q1利润缺口源于低利润率库存出货，非需求问题；低利润率AI相关库存已在Q1基本消化完毕，良率问题已解决，2Q起出货结构将转向更高利润率产品组合。
{\small\color{refgray}数据：Q1营业利润63亿日元（GS预期70亿）\par}
{\small\color{refgray}数据：全年营业利润指引从297亿上调至337亿日元\par}
{\small\color{refgray}数据：AI应用全年销售指引约310亿日元，同比翻倍以上\par}
{\small\color{refgray}数据：半导体应用需求恢复提前至Q2，订单创历史新高\par}
{\small\color{refgray}边际变化：Q1利润低于预期但全年指引上调，低利润率库存出清后出货结构将改善；AI订单确认和交付日期锁定提供可见度。\par}
\paragraph{GS} 韩国存储专家会议判断DRAM定价全年保持双位数季度环比增长，HBM定价明年上行空间可能被低估。供应持续紧张，长期协议（LTA）约束力增强（超一半服务器DRAM已覆盖），中国供应商中期难以构成实质威胁。GS预计三星电子HBM定价明年同比+87\%，高于彭博一致预期的52\%。
{\small\color{refgray}数据：DRAM Q3定价双位数环比增长，Q4可能延续\par}
{\small\color{refgray}数据：超一半服务器DRAM已签署LTA，含大额预付款和照付不议条款\par}
{\small\color{refgray}数据：GS预计三星HBM定价明年同比+87\%（彭博一致预期+52\%）\par}
{\small\color{refgray}数据：产能扩张但HBM占比提升，位元增长低于历史均值\par}
{\small\color{refgray}边际变化：专家认为HBM定价上行空间可能被市场低估，LTA覆盖率提升将结构性降低定价波动。\par}
\paragraph{GS} 国巨Q2 AI营收占比达16\%，提前突破全年15\%目标，产能利用率80-85\%且预计Q3提升至90\%以上，整体产能接近满载。订单出货比（BB ratio）从Q1的1.0+升至6月底的2.2，为下半年需求走强明确信号。客户长约需求增加，尤其集中在MLCC和钽电容。
{\small\color{refgray}数据：AI营收占比Q2达16\%（Q1为14-15\%），提前突破全年15\%目标\par}
{\small\color{refgray}数据：产能利用率Q2为80-85\%，Q3预计提升至90\%以上\par}
{\small\color{refgray}数据：BB ratio从Q1的1.0+升至6月底的2.2\par}
{\small\color{refgray}数据：高雄、墨西哥、越南、苏州等多地推进扩产\par}
{\small\color{refgray}边际变化：AI营收占比提前达标，BB ratio大幅攀升显示需求走强；客户长约需求增加，高端电容需求正在起量。\par}
\paragraph{HSBC} 意法半导体Q2订单出货比接近2，分销库存已低于公司目标水平，持续涨价且净资本支出指引上调至20-22亿美元区间上沿。Q3指引低于预期，但复苏节奏而非趋势中断。AI数据中心收入2026年预期上调至超10亿美元，2027年远超20亿美元，主要依靠硅光互联拉动。
{\small\color{refgray}数据：整体订单出货比接近2，通信和电脑外设显著高于2\par}
{\small\color{refgray}数据：分销库存已低于公司目标水平\par}
{\small\color{refgray}数据：AI数据中心收入2026年预期超10亿美元，2027年远超20亿美元\par}
{\small\color{refgray}数据：净资本支出指引上调至20-22亿美元区间上沿\par}
{\small\color{refgray}边际变化：Q3指引低于预期但订单数据较强，复苏节奏调整而非趋势中断；AI/光互联成为最大增量。\par}
\subsubsection*{关键数据}
\begin{itemize}
  \item NXP数据中心业务仅占收入3\%，中国乘用车零售同比-20\%
  \item Seagate 9月季度定价同比+25\%，增量毛利率83\%
  \item 爱德万测试全年营业利润指引从6275亿上调至8460亿日元
  \item DRAM Q3定价双位数环比增长，HBM明年GS预期同比+87\%
  \item 国巨AI营收占比Q2达16\%，BB ratio升至2.2
  \item KLA 2026年WFE预测上调至1500亿美元以上，RPO达125亿美元
  \item 意法半导体订单出货比接近2，分销库存低于目标
  \item MARUWA AI应用全年销售指引约310亿日元，同比翻倍
\end{itemize}
\begin{figure}[htbp]
\centering
\includegraphics[width=0.88\linewidth]{figures/fig_047.jpg}
\caption{Exhibit 1 - Exhibit 1: STX F4Q26 Earnings Variance}
\end{figure}
\begin{figure}[htbp]
\centering
\includegraphics[width=0.88\linewidth]{figures/fig_100.jpg}
\caption{Exhibit 1 - Exhibit 2: Yageo revenue mix by channel (2Q26) Exhibit 3: Yageo revenue mix by region (2Q26) Exhibit 4: Yageo revenue mix by product (2Q26)}
\end{figure}
\begin{figure}[htbp]
\centering
\includegraphics[width=0.88\linewidth]{figures/fig_136.jpg}
\caption{Exhibit 1 - \#\# OWNERSHIP POSITIONING Refinitiv; MSPB Content. Includes certain hedge fund exposures held with MSPB. Information may be inconsistent with or may not reflect broader market trends. Long/Short Ratio = Long Exposure / Short exposure. Sector \% of Total Net Expo}
\end{figure}
\begin{figure}[htbp]
\centering
\includegraphics[width=0.88\linewidth]{figures/fig_035.jpg}
\caption{EXHIBIT 2 - EXHIBIT 5: Automotive and Mobile were a touch above consensus, with Industrial/IoT and Comms driving the beat on consensus}
\end{figure}
{\small\color{refgray}本节综合 8 篇研究；涉及 Bernstein、MS、GS、HSBC。}
\newpage
\subsection{消费与零售：分化加剧，品牌力与渠道重塑成为核心主线}
\textbf{全球消费与零售市场呈现显著分化：奢侈品领域，爱马仕利润率超预期但区域增长引擎切换至欧洲和日本，中国高净值客群韧性犹存但大众消费承压；运动品牌中，阿迪达斯营收超预期但利润受营销成本侵蚀，Nike主动收缩中国线上批发以换取品牌质量；新茶饮与餐饮赛道，古茗开店节奏放缓但利润改善，海底捞核心业务恢复慢于预期；中国奶粉行业受质量事件冲击，头部品牌加速份额集中。}
\subsubsection*{主流共识}
\begin{itemize}
  \item 中国消费市场结构性分化加剧，高净值客群消费韧性显著强于大众消费者，奢侈品、高端运动品牌均呈现这一趋势。
  \item 品牌力与产品创新成为穿越周期的核心能力，爱马仕、阿迪达斯等通过品类组合向上延伸或功能驱动增长切换来维持竞争力。
  \item 渠道策略从追求规模转向提升质量，Nike终止中国线上批发、古茗放缓开店节奏均反映这一转变。
\end{itemize}
\subsubsection*{分歧与条件差异}
\begin{itemize}
  \item 奢侈品增长引擎的区域切换：爱马仕依赖欧洲和日本，而开云集团亚太降幅收窄但中国大陆仍在调整，市场对‘中国消费者全球消费’与‘中国境内消费’的解读存在分歧。
  \item 运动品牌利润与增长的取舍：阿迪达斯营收超预期但利润率下滑，Nike以短期收入下降换取长期品牌溢价，市场对两种策略的长期效果看法不一。
  \item 中国新茶饮与餐饮的复苏路径：古茗通过门店质量优化实现利润改善，而海底捞核心业务翻台率增长乏力、新业务稀释利润率，市场对‘旧引擎减速、新引擎未成型’的过渡期定价存在分歧。
\end{itemize}
\subsubsection*{机构观点}
\paragraph{Bernstein} 爱马仕1H26有机销售+6.7\%符合预期，但EBIT利润率41.1\%超预期70bps，毛利率提升是主因。区域分化明显：法国+6.2\%、日本+12.3\%超预期，亚太（除日本）+2.5\%略低。品类上，丝绸+12.2\%、手表+4.4\%意外走强，香水-9.5\%拖累。市场过度聚焦皮革制品产量增长（+6\%为生产工时而非产量），忽视了高级珠宝、高级定制时装等品类组合向上延伸的增长空间。账上现金超120亿欧元，特别分红概率上升。
{\small\color{refgray}数据：1H26有机销售+6.7\%，符合预期\par}
{\small\color{refgray}数据：EBIT利润率41.1\%，超预期70bps\par}
{\small\color{refgray}数据：法国Q2有机增长+6.2\%（预期+3.3\%），日本+12.3\%（预期+10.6\%）\par}
{\small\color{refgray}数据：亚太（除日本）+2.5\%（预期+3.3\%）\par}
{\small\color{refgray}边际变化：市场此前担忧皮革制品产量增速放缓，但Bernstein指出品类组合向上延伸（高级珠宝、高级定制时装）可吸收产量增长并提升利润率，这一增长来源未被充分定价。\par}
\paragraph{NOM} 爱马仕2Q26亚太（除日本）收入同比+2.5\%，低于预期+4.5\%，但管理层确认中国区销售仍在增长且出现企稳迹象。高净值客户韧性较强，而‘有抱负的消费者’面临更大困难，高价位产品（珠宝、成衣）表现更好。开云集团2Q26可比收入+2\%略超预期，Gucci跌幅从-8\%收窄至-2\%，亚太降幅从-4\%收窄至-1\%，但中国大陆仍在调整。中国消费者变得更挑剔、更注重体验、更倾向本土品牌。
{\small\color{refgray}数据：爱马仕亚太（除日本）2Q26收入+2.5\%（预期+4.5\%）\par}
{\small\color{refgray}数据：开云集团2Q26可比收入+2\%（预期36.2亿欧元）\par}
{\small\color{refgray}数据：Gucci 2Q26可比收入-2\%（1Q26为-8\%）\par}
{\small\color{refgray}数据：开云亚太（除日本）零售收入-1\%（1Q26为-4\%）\par}
{\small\color{refgray}边际变化：NOM强调中国消费增长动力从大众消费者转向高净值人群，且国际品牌需适应中国消费者更挑剔、更本土化的新偏好，这一结构性转变正在重塑奢侈品市场格局。\par}
\paragraph{Bernstein} 阿迪达斯2Q26营收货币中性+14\%超预期13\%，但营业利润率8.5\%低于预期9.4\%，主因营销费用超支1.7个百分点及DTC履约成本上升。运动表现品类+39\%成为增长引擎（足球、跑步、赛车驱动），生活方式仅+2\%。北美+17\%、大中华+15\%表现强劲，欧洲批发保守仅+6\%。Nike中国线上渠道调整：2027年1月起终止合作伙伴线上店铺，预计带来约10亿美元收入下滑但毛利率改善约200bps。Nike中国份额从2020年27\%降至2025年约16\%。
{\small\color{refgray}数据：阿迪达斯2Q26营收货币中性+14\%（预期+13\%）\par}
{\small\color{refgray}数据：营业利润率8.5\%（预期9.4\%，去年同期9.2\%）\par}
{\small\color{refgray}数据：运动表现品类+39\%，生活方式仅+2\%\par}
{\small\color{refgray}数据：Nike中国线上批发调整影响约10亿美元收入\par}
{\small\color{refgray}边际变化：阿迪达斯增长引擎从潮流驱动切换至功能驱动，但利润端受营销和履约成本拖累；Nike主动收缩中国线上批发，以短期收入下降换取品牌质量和定价能力，这是其渠道策略的重大转向。\par}
\paragraph{GS} 古茗控股1H26预计营收+30\%、调整后净利润+43\%，但2026年净新增门店预期从3600家降至2000家，主因加盟商意愿减弱和公司优先门店质量。全年单店GMV预计同比-2\%，但毛利率同比+1.6个百分点。海底捞1H26收入预计+8\%至223.55亿元，净利润+5\%至18.47亿元，翻台率增长约2\%但客单价下滑约1\%，全年单店收入预期从+3\%下调至-0.5\%。新业务增长快但稀释利润率。
{\small\color{refgray}数据：古茗1H26预计营收+30\%至73.72亿元\par}
{\small\color{refgray}数据：调整后净利润+43\%至16.30亿元\par}
{\small\color{refgray}数据：2026年净新增门店预期从3600家降至2000家\par}
{\small\color{refgray}数据：海底捞1H26收入预计+8\%至223.55亿元\par}
{\small\color{refgray}边际变化：GS下调古茗和海底捞盈利预测，核心变化在于：古茗从追求开店数量转向提升单店质量，海底捞核心业务从‘量价齐升’转向‘量稳价降’，新业务尚不能弥补利润缺口。\par}
\paragraph{GS} 中国婴幼儿配方奶粉行业7月连续出现两起质量事件：Cabio因ARA问题发布盈利预警，预计1H26净亏损约1.01亿元，触发ST风险警示；Friso澳门批次铅含量超标，但第三方复检结果远低于限值。行业线上销售已连续两个季度下滑，但头部品牌加速份额集中，贝拉米和Biostime 6月线上销售额同比分别+70\%和+43\%。
{\small\color{refgray}数据：Cabio预计1H26净亏损约1.01亿元\par}
{\small\color{refgray}数据：Friso澳门批次铅含量0.02mg/kg（复检<0.007mg/kg）\par}
{\small\color{refgray}数据：贝拉米6月线上销售额同比+70\%\par}
{\small\color{refgray}数据：Biostime 6月线上销售额同比+43\%\par}
{\small\color{refgray}边际变化：奶粉行业从需求调整、渠道变化转向质量事件冲击，消费者信心短期承压，但强品质管控的头部品牌（贝拉米、Biostime）反而加速扩大份额，行业集中度提升趋势强化。\par}
\paragraph{JPM} 截至7月29日当周，散户资金流降至过去一年历史百分位的0.8\%，连续四天净卖出个股。科技硬件股遭遇有记录以来最大单周资金流出，半导体板块拥挤度从99百分位降至91.8百分位。科技ETF流出规模为历史第二，主要驱动力来自三倍做多半导体ETF（SOXL）。美光从散户最青睐名单跌入最卖名单。
{\small\color{refgray}数据：散户资金流降至历史百分位0.8\%\par}
{\small\color{refgray}数据：科技硬件股单周流出创纪录\par}
{\small\color{refgray}数据：半导体拥挤度从99\%降至91.8\%\par}
{\small\color{refgray}数据：SOXL累计零售净买入额出现拐点\par}
{\small\color{refgray}边际变化：JPM指出本轮减仓并非简单避险，而是动量策略反转叠加持仓变化后的被动调整，散户集中清理此前拥挤度最高的方向（半导体、硬件），而非全面离场。\par}
\subsubsection*{关键数据}
\begin{itemize}
  \item 爱马仕1H26 EBIT利润率41.1\%，超预期70bps
  \item 阿迪达斯2Q26运动表现品类+39\%，生活方式仅+2\%
  \item Nike中国份额从27\%降至16\%，线上批发调整影响约10亿美元收入
  \item 古茗2026年净新增门店预期从3600家降至2000家
  \item 海底捞全年单店收入预期从+3\%下调至-0.5\%
  \item Cabio预计1H26净亏损约1.01亿元，触发ST风险警示
  \item 贝拉米6月线上销售额同比+70\%
  \item 散户资金流降至历史百分位0.8\%，半导体拥挤度从99\%降至91.8\%
\end{itemize}
\begin{figure}[htbp]
\centering
\includegraphics[width=0.88\linewidth]{figures/fig_022.jpg}
\caption{EXHIBIT 3 - EXHIBIT 3: Luxury Goods Calendar Year 2Q26 results CY 2Q26 Reported Organic Growth and Actual Change in Revenue \#\# INVESTMENT IMPLICATIONS \#\# We rate Hermès Outperform with a target price of EUR2,150.}
\end{figure}
\begin{figure}[htbp]
\centering
\includegraphics[width=0.88\linewidth]{figures/fig_064.jpg}
\caption{EXHIBIT 1 - EXHIBIT 1: Nike's share has declined from around 27\% at its peak to around 16\% in 2025, with much of the pre-COVID share expansion now having been reversed EXHIBIT 2: Since 2020, local competitors have steadily taken share from I}
\end{figure}
\begin{figure}[htbp]
\centering
\includegraphics[width=0.88\linewidth]{figures/fig_087.jpg}
\caption{Exhibit 1 - Exhibit 1: MNC players saw sequential market share trend mixed in Mar-Apr, Friesland/Danone/A2 value share increased while Wyeth decreased vs. Jan-Feb Offline international IMF brands market share by value Exhibit 2: IMF 618 Shop}
\end{figure}
\begin{figure}[htbp]
\centering
\includegraphics[width=0.88\linewidth]{figures/fig_116.jpg}
\caption{Figure 54 - Figure 1: Momentum Crowding Figure 2: Momentum Unwind Now vs Feb-Mar'25 Figure 3: Crowding in Low Quality}
\end{figure}
{\small\color{refgray}本节综合 10 篇研究；涉及 Bernstein、NOM、GS、JPM。}
\newpage
\subsection{AI变现进入验证期：从广告效率到企业级代理}
\textbf{2026年Q2，全球互联网与传媒板块的核心叙事从AI投入转向AI变现验证。Meta与微软分别通过广告效率提升和云/AI代理订阅证明了AI投入的短期回报，但长期资本支出路径与消费者AI战略的不确定性仍是市场分歧焦点。中国互联网则呈现流量向企业级AI代理平台转移的结构性变化。}
\subsubsection*{主流共识}
\begin{itemize}
  \item AI投入在核心广告业务中已产生可量化回报：Meta广告收入同比增长27-28\%，Advantage+年化收入超750亿美元，AI驱动的推荐系统显著提升点击率与转化率。
  \item 微软Azure增速连续加速至43\%，M365 Copilot付费席位单季净增1000万至3000万以上，验证了企业级AI订阅的强劲需求。
  \item 中国互联网流量正从消费级AI聊天机器人向企业级AI代理平台转移，腾讯元宝月访问量从885万激增至2100万。
\end{itemize}
\subsubsection*{分歧与条件差异}
\begin{itemize}
  \item Meta消费者AI战略（个人超级智能）与谷歌/OpenAI的企业级AI路径存在根本分歧，前者依赖36亿DAU的消费者生态，后者聚焦企业变现。
  \item Meta管理层对2027-2028年资本支出路径缺乏具体框架，GS上调累计资本支出预测约14\%，而市场对长期自由现金流轨迹存在疑虑。
  \item 中国AI聊天机器人领域，豆包DAU达1.67亿且粘性高，通义千问DAU虽达2930万但55\%为轻度用户，用户质量差异显著。
\end{itemize}
\subsubsection*{机构观点}
\paragraph{NOM} 中国移动互联网6月总使用时长同比增长8.6\%，流量正从消费级AI聊天机器人向企业级AI代理平台转移。字节豆包DAU达1.67亿，环比增6\%，同比增3.3倍，继续扩大领先；DeepSeek以3240万DAU居次。通义千问DAU达2930万，但55\%为轻度用户（日均使用不足1分钟），粘性远低于豆包（85\%重度用户）和DeepSeek（81\%）。企业AI代理平台中，腾讯元宝月访问量从3月885万升至2100万，字节TRAE IDE以1279万次居次。
{\small\color{refgray}数据：中国6月活跃智能手机用户12.8亿，同比增1.2\%\par}
{\small\color{refgray}数据：总使用时长同比增8.6\%，较5月的10.1\%放缓\par}
{\small\color{refgray}数据：豆包DAU 1.67亿，环比增6\%，同比增3.3倍\par}
{\small\color{refgray}数据：通义千问DAU 2930万，55\%为轻度用户\par}
{\small\color{refgray}边际变化：流量从消费级AI聊天机器人向企业级AI代理平台转移，通义千问用户增长快但粘性恶化，轻度用户占比从28\%升至55\%。\par}
\paragraph{Bernstein} Meta 2Q26营收增长28\%，广告收入增27\%至594亿美元，展示量增14\%，定价增12\%。AI驱动的Advantage+年化收入超750亿美元，生成式系统推动Instagram应用事件转化率提升1\%，Facebook广告点击增8.3\%、转化增15.7\%。但核心关注点是扎克伯格提出的“个人超级智能”消费者AI战略，与谷歌/OpenAI的企业级路径形成分歧。Meta押注消费者AI，依赖36亿DAU生态，旨在避免下一个“TikTok式”冲击。
{\small\color{refgray}数据：Meta 2Q26广告收入594亿美元，同比增27\%\par}
{\small\color{refgray}数据：Advantage+年化收入超750亿美元\par}
{\small\color{refgray}数据：Facebook广告点击增8.3\%，转化增15.7\%\par}
{\small\color{refgray}数据：Instagram应用事件转化率提升1\%\par}
{\small\color{refgray}边际变化：Meta消费者AI战略从概念走向具体，扎克伯格年度备忘录明确“个人超级智能”愿景，与同行企业级AI路径分化。\par}
\paragraph{GS} Meta 2Q26广告收入同比增28\%，AI驱动的Advantage+年化收入达750亿美元，内容推荐改进和变现效率提升是主要驱动力。管理层将AI代理视为核心业务之外的延伸，潜在机会包括个性化消费者AI代理、广告主/创作者AI工具，但变现路径仍不清晰。企业级AI机会涵盖基础模型API、业务构建工具和计算直接变现，但规模与路径尚不明确。GS将2027-2028年累计资本支出预测上调约14\%，反映管理层“短期优先建设计算能力”的表述。
{\small\color{refgray}数据：Meta 2Q26广告收入同比增28\%\par}
{\small\color{refgray}数据：Advantage+年化收入750亿美元\par}
{\small\color{refgray}数据：GS上调2027-2028年累计资本支出预测约14\%\par}
{\small\color{refgray}数据：管理层短期优先建设计算能力，资本支出继续上升\par}
{\small\color{refgray}边际变化：Meta从定性描述AI投入转向提供短期回报证据，但长期资本支出框架仍缺失，GS上调资本支出预测反映市场对自由现金流轨迹的疑虑。\par}
\paragraph{MS} 微软F4Q26 Azure固定汇率增长43\%，超指引上限3个百分点，F1Q27指引约45\%，增长瓶颈在供给端。M365 Copilot付费席位单季净增约1000万至3000万以上，超市场预期；5万+席位大客户同比增7倍，全员部署客户环比增75\%。变现模式从单一座位许可向三引擎ARPU框架迁移：直接Copilot席位、E7订阅、消费型变现。运营利润率45.1\%仍同比扩张，AI投入未明显侵蚀利润。Meta核心社交业务日活36亿，Instagram突破20亿DAU，AI+LLM提升广告匹配效率，3Q26收入指引上限640亿美元隐含同比增26\%。四项新产品期权（Neocloud、API、订阅、消息商业化）合计可为2028年EPS贡献约9美元上行空间。
{\small\color{refgray}数据：Azure F4Q26固定汇率增长43\%，超指引上限3个百分点\par}
{\small\color{refgray}数据：M365 Copilot付费席位突破3000万，单季净增1000万\par}
{\small\color{refgray}数据：运营利润率45.1\%，超预期70个基点\par}
{\small\color{refgray}数据：Meta 3Q26收入指引上限640亿美元，同比增26\%\par}
{\small\color{refgray}边际变化：微软AI变现从“能否发生”进入“能持续多久”阶段，Azure供给瓶颈和Copilot采用加速是核心变化；Meta新产品期权从概念进入“soon”阶段，Connect大会成为关键节点。\par}
\subsubsection*{关键数据}
\begin{itemize}
  \item Meta 2Q26广告收入594亿美元，同比增27\%
  \item Advantage+年化收入超750亿美元
  \item Azure F4Q26固定汇率增长43\%，超指引上限3个百分点
  \item M365 Copilot付费席位突破3000万，单季净增1000万
  \item Meta日活跃用户36亿，Instagram突破20亿DAU
  \item 豆包DAU 1.67亿，同比增3.3倍
  \item 腾讯元宝月访问量2100万，较3月增137\%
  \item GS上调Meta 2027-2028年累计资本支出预测约14\%
\end{itemize}
\begin{figure}[htbp]
\centering
\includegraphics[width=0.88\linewidth]{figures/fig_059.jpg}
\caption{EXHIBIT 2 - EXHIBIT 2: Meta 2Q26: EBIT is impacted by one-time legal charge EXHIBIT 3: Digital ads growth (\%) Reported ad revenue Y/Y growth (\%) through 2Q26 EXHIBIT 4: Revenue contribution by geography}
\end{figure}
\begin{figure}[htbp]
\centering
\includegraphics[width=0.88\linewidth]{figures/fig_137.jpg}
\caption{Exhibit 1 - Exhibit 1: Model Changes}
\end{figure}
\begin{figure}[htbp]
\centering
\includegraphics[width=0.88\linewidth]{figures/fig_138.jpg}
\caption{Exhibit 1 - Exhibit 2: Our \textbackslash{}\$775 PT Implies Paying 23X our average '27/'28 EPS of \textbackslash{}\$34/\textbackslash{}\$35, representing 40\% upside}
\end{figure}
\begin{figure}[htbp]
\centering
\includegraphics[width=0.88\linewidth]{figures/fig_060.jpg}
\caption{EXHIBIT 4 - EXHIBIT 4: Revenue contribution by geography EXHIBIT 5: Y/Y Revenue growth by geography EXHIBIT 6: Ad impressions vs. pricing}
\end{figure}
{\small\color{refgray}本节综合 5 篇研究；涉及 NOM、Bernstein、GS、MS。}
\newpage
\subsection{工业与能源：电网投资与AI驱动的结构性增长}
\textbf{全球电网投资加速与AI数据中心需求爆发是当前工业与能源板块的核心驱动力，但区域分化、政策风险与供应链瓶颈带来不确定性。中国变压器出口结构性增长、高端CCL供需趋紧、工业自动化龙头业绩超预期，共同指向结构性机会，而美国逆变器政策与SpaceX解禁压力构成短期风险。}
\subsubsection*{主流共识}
\begin{itemize}
  \item 全球电网投资持续高景气，变压器等关键设备供需紧张，中国供应商凭借供应链优势填补缺口。
  \item AI数据中心需求驱动高端CCL、电力设备等板块量价齐升，相关企业毛利率进入扩张周期。
  \item 工业自动化领域，中国市场需求回暖，施耐德、西门子等龙头业绩超预期，行业资本开支周期或已触底。
\end{itemize}
\subsubsection*{分歧与条件差异}
\begin{itemize}
  \item 美国电力变压器缺口收窄速度：GS预计从72\%收窄至2028年的57\%，但部分机构认为本土产能扩张可能快于预期。
  \item 高端CCL产能锁定程度：GS认为2028年前产能已被提前预订，但市场对AI需求持续性存疑。
  \item 美国逆变器政策影响：NOM认为第三国产能无法规避限制，但部分市场参与者认为东南亚产能仍有机会。
  \item SpaceX解禁压力：MS认为技术面压力大于基本面，但多空分歧明显，部分投资者认为估值已反映风险。
\end{itemize}
\subsubsection*{机构观点}
\paragraph{Bernstein} 霍尼韦尔二季度有机收入同比+4\%（预期0\%），EPS \$1.95超预期7\%，三大业务板块均超预期。管理层下半年增长指引4-6\%偏保守，订单增速16\%且订单出货比>1.1倍提供支撑。工业自动化利润率从17\%向22\%扩张路径清晰，一半来自剥离低毛利业务，一半来自经营杠杆。市场对管理层可信度的评估正在改善。
{\small\color{refgray}数据：有机收入同比+4\%，市场预期0\%\par}
{\small\color{refgray}数据：调整后EPS \$1.95，超共识7\%\par}
{\small\color{refgray}数据：总订单有机增长16\%，各板块两位数增长\par}
{\small\color{refgray}数据：订单出货比>1.1倍\par}
{\small\color{refgray}边际变化：此前持怀疑态度的分析师开始重新评估，二季度数据使'展示给我看'故事变得可信。\par}
\paragraph{MS} SpaceX二季度收入预计\$67.5亿略低于共识，但真正压力来自8月6日首批9.3亿股解禁（约\$1000亿市值），总计40亿股将在2027年1月前解锁。报价从高点回落近50\%，低于IPO发行价。超过2/3投资者对年底前表现持谨慎态度，空头信心增强。管理层评论比季度数字更重要，首个明确催化剂是8月底Starship第14次试飞。
{\small\color{refgray}数据：二季度收入预计\$67.5亿，共识\$69.2亿\par}
{\small\color{refgray}数据：8月6日首批9.3亿股解禁，约\$1000亿市值\par}
{\small\color{refgray}数据：总计40亿股在2027年1月前解锁\par}
{\small\color{refgray}数据：报价从高点回落近50\%\par}
{\small\color{refgray}边际变化：焦点从基本面转向技术面解禁压力，市场情绪从乐观转为谨慎。\par}
\paragraph{NOM} FCC将外国产逆变器纳入覆盖清单，阳光电源产品组合大概率受限。标准采用《购买美国产品法》国内最终产品测试，第三国产能无法合规。短期祖父条款提供缓冲至2026年，但新机型认证通道关闭将侵蚀竞争力。固态变压器等第二增长曲线可能被削弱，硬件去中国化+软件本地化推高成本。下调2027/28年EPS至8.78/10.41元。
{\small\color{refgray}数据：FCC覆盖清单明确包含混合电池逆变器\par}
{\small\color{refgray}数据：采用Buy America测试，非原产地规则\par}
{\small\color{refgray}数据：2027/28年EPS下调至8.78/10.41元（原9.13/11.20元）\par}
{\small\color{refgray}数据：目标价从16倍2026年PE调整至14倍2027年PE\par}
{\small\color{refgray}边际变化：市场此前认为东南亚产能可规避限制，但FCC援引Buy America测试使第三国产能策略根本性存疑。\par}
\paragraph{GS} 中国变压器6月出口同比+47\%（5月+72\%），但对美出口同比+172\%、环比+52\%。美国本土供需缺口当前72\%，预计2028年收窄至57\%，全球交货周期128周，中国供应商仅24-36周。中东（+141\%）和美洲（+108\%）是主要增量区域。220-330MVA段出口均价近3月滚动同比-8\%，但6月单月已转正为+3\%。
{\small\color{refgray}数据：6月总出口同比+47\%，对美出口+172\%\par}
{\small\color{refgray}数据：美国供需缺口72\%，预计2028年收窄至57\%\par}
{\small\color{refgray}数据：全球交货周期128周，中国供应商24-36周\par}
{\small\color{refgray}数据：中东出口+141\%，美洲+108\%\par}
{\small\color{refgray}边际变化：对美出口增速（+172\%）与整体增速（+47\%）背离扩大，美国供需紧张未缓解，中国供应商持续填补缺口。\par}
\paragraph{GS} 高端CCL供需从'产能紧张'转向'产能被提前锁定'。台光2Q26毛利率33.9\%（环比+4.4ppt），台燿29.7\%（+4.6ppt），均超预期。3Q26展望：台光营收环比+21\%，毛利率升至36.8\%；台燿营收环比+23\%，毛利率再扩3.4ppt。Trainium 3/Rubin等新AI项目驱动M8/M9级CCL需求爬坡，全产品线提价10-60\%。2028年前产能已被预订。
{\small\color{refgray}数据：台光2Q26毛利率33.9\%，环比+4.4ppt\par}
{\small\color{refgray}数据：台燿2Q26毛利率29.7\%，环比+4.6ppt\par}
{\small\color{refgray}数据：台光3Q26预计毛利率36.8\%\par}
{\small\color{refgray}数据：全产品线提价10-60\%\par}
{\small\color{refgray}边际变化：毛利率扩张驱动力从量价齐升切换为定价权提升与产品结构改善的双轮驱动，基板级CCL业务即将贡献结构性支撑。\par}
\paragraph{GS} 宏发股份2Q26收入59.14亿元同比+36\%，毛利率33\%。2026年预计收入增长27\%，远高于过去十年15\%的CAGR，其中约10个百分点来自提价。毛利率预计3Q26回升至34.5\%、4Q26至35.1\%。六大产品线均两位数增长，HVDC继电器+30\%，通用继电器+36\%。欧姆龙剥离继电器业务为宏发提供份额扩张窗口。
{\small\color{refgray}数据：2Q26收入59.14亿元，同比+36\%\par}
{\small\color{refgray}数据：2026年预计收入增长27\%，过去十年CAGR 15\%\par}
{\small\color{refgray}数据：价格因素贡献>10个百分点收入增长\par}
{\small\color{refgray}数据：毛利率预计3Q26回升至34.5\%\par}
{\small\color{refgray}边际变化：增长逻辑从跟随行业转向主动夺取份额，欧姆龙业务剥离标志竞争格局根本性变化。\par}
\paragraph{GS} 日立第一财季能源业务订单同比+87\%（剔除核电+94\%），订单积压10.3万亿日元环比+12\%。调整后EBITA 1290亿日元超预期，公司上调全年指引。基础订单（变压器、开关设备）同样强劲，积压订单利润率改善。AI提效目标：FY3/26末生产率提升约10\%，目标FY3/28达30\%。上调FY3/27-29盈利预测6\%/5\%/4\%。
{\small\color{refgray}数据：能源订单同比+87\%，剔除核电+94\%\par}
{\small\color{refgray}数据：订单积压10.3万亿日元，环比+12\%\par}
{\small\color{refgray}数据：AI提效目标：FY3/28达30\%\par}
{\small\color{refgray}数据：上调FY3/27-29盈利预测6\%/5\%/4\%\par}
{\small\color{refgray}边际变化：能源业务订单增长并非仅靠大单拉动，基础订单广泛走强，积压订单利润率改善为未来利润释放提供可见度。\par}
\paragraph{GS} 力拓1H26基础EBITDA \$148亿同比+28\%，净利润\$69亿同比+43\%超预期11\%。有效税率25\%低于指引30\%。22个主要开发项目全部按计划或超前推进，成本削减目标从\$6.5亿提升至\$18亿（约占成本基础5\%），上半年已实现\$13亿年化。非核心资产出售计划\$50-100亿，预计2026年底前释放约\$50亿。
{\small\color{refgray}数据：1H26 EBITDA \$148亿，同比+28\%\par}
{\small\color{refgray}数据：净利润\$69亿，同比+43\%\par}
{\small\color{refgray}数据：有效税率25\%，低于指引30\%\par}
{\small\color{refgray}数据：成本削减目标\$18亿，占成本基础5\%\par}
{\small\color{refgray}边际变化：成本削减目标从\$6.5亿提升至\$18亿，执行力比市场预期更激进，资产组合盈利结构正在改变。\par}
\paragraph{GS} 施耐德电气2Q26有机销售增长16.5\%，超预期近6个百分点。上调全年指引：有机增长从7-10\%至10-13\%，利润率从19.1-19.4\%至19.4-19.7\%。中国工业自动化增速19.6\%，远超预期7.7\%，与西门子信号形成正向共振。数据中心三位数增长，住宅业务在欧洲环比改善。调整后EBITA利润率16.5\%，高于ABB的12\%。
{\small\color{refgray}数据：2Q26有机销售增长16.5\%，预期10.6\%\par}
{\small\color{refgray}数据：全年指引上调至10-13\%\par}
{\small\color{refgray}数据：中国工业自动化增速19.6\%，预期7.7\%\par}
{\small\color{refgray}数据：调整后EBITA利润率16.5\%\par}
{\small\color{refgray}边际变化：中国工业自动化19.6\%增速是最大亮点，可能预示中国制造业资本开支边际回暖，与西门子信号形成共振。\par}
\paragraph{MS} Bloom Energy 2Q26营收\$10.65亿首次破10亿，超预期29\%。连续第二个季度上调全年指引：营收中位数上调约\$4.5亿，利润指引上调26\%。调整后毛利率34.3\%，超预期350bps。经营现金流\$2.26亿，市场预期为负。管理层对钪供应、项目延迟、产能瓶颈三大担忧均给出增量安抚，backlog增速超收入增速，推算设备backlog至少\$80亿。
{\small\color{refgray}数据：2Q26营收\$10.65亿，超预期29\%\par}
{\small\color{refgray}数据：全年营收指引上调约\$4.5亿\par}
{\small\color{refgray}数据：调整后毛利率34.3\%，超预期350bps\par}
{\small\color{refgray}数据：经营现金流\$2.26亿，市场预期-0.42亿\par}
{\small\color{refgray}边际变化：连续两个季度主动上调指引，管理层对订单可见度和交付节奏信心增强，运营杠杆开始显现。\par}
\subsubsection*{关键数据}
\begin{itemize}
  \item 中国变压器6月对美出口同比+172\%，整体+47\%
  \item 美国变压器供需缺口72\%，预计2028年收窄至57\%
  \item 高端CCL台光3Q26毛利率预计36.8\%
  \item 施耐德中国工业自动化增速19.6\%，远超预期7.7\%
  \item 宏发股份2026年收入增长27\%，价格贡献>10ppt
  \item 日立能源订单同比+87\%，积压10.3万亿日元
  \item 力拓成本削减目标从\$6.5亿提升至\$18亿
  \item Bloom Energy 2Q26营收\$10.65亿，经营现金流\$2.26亿
\end{itemize}
\begin{figure}[htbp]
\centering
\includegraphics[width=0.88\linewidth]{figures/fig_082.jpg}
\caption{Exhibit 1 - Exhibit 2: China total transformer export value picked up to \$47\textbackslash{}\%\$ yoy in June (vs \$72\textbackslash{}\%\$ in May)}
\end{figure}
\begin{figure}[htbp]
\centering
\includegraphics[width=0.88\linewidth]{figures/fig_090.jpg}
\caption{Exhibit 2 - Exhibit 3: We expect both EMC and TUC GM to start the next round of expansion driven by solid AI demand and tightening S/D outlook Exhibit 4: Both EMC and TUC will expand more than \$100\textbackslash{}\%+\$ capacity in coming 2 years, while both}
\end{figure}
\begin{figure}[htbp]
\centering
\includegraphics[width=0.88\linewidth]{figures/fig_104.jpg}
\caption{Exhibit 1 - Exhibit 1: Schneider 2Q26 results vs. company-compiled consensus Exhibit 2: Schneider Group OSG vs peers Exhibit 3: Schneider Energy Management OSG vs peers We are Buy rated on Schneider (also on the Conviction List), with a 12}
\end{figure}
\begin{figure}[htbp]
\centering
\includegraphics[width=0.88\linewidth]{figures/fig_097.jpg}
\caption{Exhibit 2 - Exhibit 2: PT, NAV and earnings revisions}
\end{figure}
{\small\color{refgray}本节综合 10 篇研究；涉及 Bernstein、MS、NOM、GS。}
\newpage
\subsection{汽车与出行：美国定价稳定与产品周期分化}
\textbf{美国市场定价环境保持稳定，为高敞口日系车企提供盈利支撑；中国乘用车市场新车发布效应消退，订单分化加剧，行业等待下一轮产品周期。}
\subsubsection*{主流共识}
\begin{itemize}
  \item 美国市场定价环境稳定，福特上调全年EBIT指引，皮卡需求强劲，V8发动机占比上升，全年定价预计同比增长约0.5\%。
  \item 日本车企（丰田、本田、斯巴鲁）因高美国敞口和强混动组合，受益于稳定的美国定价环境。
  \item 中国乘用车市场7月第三周多数车企订单环比回落，新车发布效应正在消退，市场进入过渡阶段。
\end{itemize}
\subsubsection*{分歧与条件差异}
\begin{itemize}
  \item 保时捷盈利质量改善（单车收入12.6万欧元，高于预期）与现金流超预期（净现金流利润率6.5\%），但市场对其能否持续回到10\%以上利润率存在分歧。
  \item 中国车企中，理想和小鹏因新车型（L6、MONA M03）订单环比仍正增长，而比亚迪、吉利等处于产品空窗期，订单表现偏弱。
  \item 福特储能系统（ESS）订单积压增长，目标2028年20GWh产能，但市场对日本车企（如本田）在美国ESS业务的发展潜力看法不一。
\end{itemize}
\subsubsection*{机构观点}
\paragraph{GS} 福特Q2财报显示美国定价环境稳定，皮卡需求强劲，全年定价预计同比增长0.5\%，这对高美国敞口的日本车企（丰田、本田、斯巴鲁）是积极信号。同时，福特储能系统订单积压增长，目标2028年20GWh产能，可能推动市场重新关注本田在美国的ESS业务。维持对这三家日本车企的积极看法。
{\small\color{refgray}数据：福特全年调整后EBIT指引从85-105亿美元上调至100-110亿美元\par}
{\small\color{refgray}数据：Ford Blue全年EBIT展望从45-50亿美元上调至50-55亿美元\par}
{\small\color{refgray}数据：美国市场全年定价预计同比增长约0.5\%\par}
{\small\color{refgray}数据：福特储能系统目标2028年实现20GWh产能\par}
{\small\color{refgray}边际变化：此前市场担忧美国定价竞争，福特财报显示定价环境稳定且向上，储能需求成为新关注点。\par}
\paragraph{GS} 保时捷2Q26业绩超预期，集团EBIT 7.53亿欧元（高于市场共识7.13亿欧元），营业利润率8.5\%（高于市场共识8.1\%），主要受911高配车型（GTS、Turbo、GT）驱动，单车收入达12.6万欧元（高于预期12.5万欧元）。汽车业务净现金流5.06亿欧元，高于GS预期3.44亿欧元，净现金流利润率6.5\%超过全年3\%-5\%指引上限。全年指引保持不变，但2027年仍有重组支出。
{\small\color{refgray}数据：保时捷2Q26集团EBIT 7.53亿欧元，高于市场共识7.13亿欧元\par}
{\small\color{refgray}数据：营业利润率8.5\%，高于市场共识8.1\%\par}
{\small\color{refgray}数据：单车收入12.6万欧元，高于GS预期12.5万欧元\par}
{\small\color{refgray}数据：汽车业务净现金流5.06亿欧元，高于GS预期3.44亿欧元\par}
{\small\color{refgray}边际变化：此前市场关注保时捷利润率压力，本次盈利质量改善（高配车型驱动、现金流超预期）显示运营韧性，但税率波动（实际税率17\% vs GS模型30\%）放大了单季EPS超预期。\par}
\paragraph{MS} 中国乘用车市场7月第三周（20-26日）多数车企订单环比回落，新车发布效应消退。比亚迪周订单7.05-7.1万辆（环比-8\%，同比-7\%），吉利银河、蔚来、零跑等均环比下滑。例外是理想（周订单9.2-9.4万辆，7月累计环比+3\%，同比+18\%）和小鹏（周订单9.9-10.1万辆，7月累计环比+43\%，同比+13\%），受益于新车型L6和MONA M03。市场等待下一轮产品周期，关键催化剂来自比亚迪的“大汉”和吉利的新车型。
{\small\color{refgray}数据：比亚迪周订单7.05-7.1万辆，环比-8\%，同比-7\%\par}
{\small\color{refgray}数据：理想周订单9.2-9.4万辆，7月累计环比+3\%，同比+18\%\par}
{\small\color{refgray}数据：小鹏周订单9.9-10.1万辆，7月累计环比+43\%，同比+13\%\par}
{\small\color{refgray}数据：蔚来周订单1.15-1.2万辆，环比-3\%，同比+121\%\par}
{\small\color{refgray}边际变化：此前7月前几周新车发布带来订单脉冲，现多数车企订单正常化回落，理想和小鹏因新车型仍保持月度环比正增长，市场焦点转向下一轮产品周期。\par}
\subsubsection*{关键数据}
\begin{itemize}
  \item 福特全年调整后EBIT指引从85-105亿美元上调至100-110亿美元
  \item 美国市场全年定价预计同比增长约0.5\%
  \item 保时捷2Q26单车收入12.6万欧元，高于预期12.5万欧元
  \item 保时捷汽车业务净现金流利润率6.5\%，超过全年3\%-5\%指引上限
  \item 比亚迪周订单7.05-7.1万辆，环比-8\%，同比-7\%
  \item 理想7月累计订单环比+3\%，同比+18\%
  \item 小鹏7月累计订单环比+43\%，同比+13\%
  \item 特斯拉中国周订单同比-32\%
\end{itemize}
{\small\color{refgray}本节综合 3 篇研究；涉及 GS、MS。}
\newpage
\subsection{医疗与保险：分化加速，创新与资本效率成关键}
\textbf{中国医疗与保险市场正经历结构性分化：跨国药企在华业绩因产品组合与政策暴露度不同而显著分化，保险行业则因资本效率差异加速两极分化。创新产品准入与成本控制能力成为区分赢家与输家的核心变量。}
\subsubsection*{主流共识}
\begin{itemize}
  \item 跨国药企在华业务受集采和医保定价改革影响，业绩表现分化明显，创新产品成为增长关键。
  \item 中国保险市场渗透率缓慢提升，但行业集中度持续上升，头部公司占据主导地位。
  \item 中小险企面临资本补充压力，核心偿付能力低于100\%的公司数量增加，市场出清加速。
  \item GSK等药企通过成本节约计划为研发提供资金，但市场已部分预期其利润率改善。
\end{itemize}
\subsubsection*{分歧与条件差异}
\begin{itemize}
  \item 跨国药企在华业绩：诺华因Leqvio纳入医保实现12\%增长，而罗氏药品销售下滑33\%，显示医保准入效果因产品而异。
  \item 保险行业资本效率：上市险企核心偿付能力普遍高于100\%，而17家未上市险企低于100\%，其中4家低于70\%，资本补充需求紧迫性存在差异。
  \item GSK成本节约的增量贡献：JPM测算额外成本节约仅带来约20个基点利润率提升，但市场共识已包含40个基点，分歧在于成本节约能否转化为管线价值。
  \item CDMO领域活跃度：龙沙等CDMO企业持续活跃，但跨国药企自建产能与外包策略的平衡点存在不确定性。
\end{itemize}
\subsubsection*{机构观点}
\paragraph{NOM} 跨国药企在华业绩分化显著：诺华中国区销售同比增长12\%，Leqvio纳入医保后使用量快速提升，部分抵消Cosentyx下滑；阿斯利康中国区销售下降7\%，受集采持续影响，但管理层对全年增长保持信心；罗氏中国区药品销售同比下滑33\%，诊断业务下滑15\%，医疗定价改革影响仍在，但管理层预期下半年减弱。创新产品准入与政策暴露度是分化核心。
{\small\color{refgray}数据：诺华中国区销售同比增长12\%（剔除汇率影响后增长6\%）至12亿美元\par}
{\small\color{refgray}数据：阿斯利康中国区销售同比下降7\%（剔除汇率影响后下降13\%）至15.9亿美元\par}
{\small\color{refgray}数据：罗氏中国区药品销售同比下滑33\%至5.22亿瑞士法郎\par}
{\small\color{refgray}数据：罗氏诊断业务同比下滑15\%至3.9亿瑞士法郎\par}
{\small\color{refgray}边际变化：此前市场普遍认为跨国药企在华业务受政策影响一致，但二季度数据显示分化加剧，诺华因创新药医保准入实现增长，而罗氏因诊断业务定价改革承压。\par}
\paragraph{JPM} 中国保险市场两极分化加速：头部五家寿险和非寿险公司分别占据45\%和68\%市场份额，而中小险企资本缓冲不足。截至2026年3月，17家未上市寿险公司核心偿付能力低于100\%，其中4家低于70\%，逼近50\%监管红线。2025年55\%的财险公司出现承保亏损。规模成为区分竞争力的关键变量，中小险企在品牌、渠道和资本方面均处劣势，预计将逐步让出市场份额。
{\small\color{refgray}数据：前5大寿险公司市场份额45\%，前5大非寿险公司市场份额68\%\par}
{\small\color{refgray}数据：17家未上市寿险公司核心偿付能力低于100\%，其中4家低于70\%\par}
{\small\color{refgray}数据：2025年55\%的财险公司出现承保亏损\par}
{\small\color{refgray}数据：保险深度从2015年3.5\%升至2025年4.4\%\par}
{\small\color{refgray}边际变化：此前市场关注保费增速放缓，但当前焦点转向资本效率与偿付能力分化，中小险企资本补充需求紧迫性上升，行业出清预期增强。\par}
\paragraph{JPM} GSK重申2031年超400亿英镑销售目标，并宣布到2029年实现19亿英镑年化成本节约的重组计划。JPM测算成本节约对核心EBIT净贡献约2亿英镑（约60个基点利润率提升），但市场共识已包含2026-2030年间约40个基点利润率增长，额外成本节约仅带来约20个基点增量，对应2030年核心EBIT约1\%上调空间。关键三期临床数据集中在2028-2030年读出，真正驱动预测上调需等待管线进展。
{\small\color{refgray}数据：GSK 2031年销售目标超400亿英镑\par}
{\small\color{refgray}数据：重组计划到2029年实现19亿英镑年化成本节约\par}
{\small\color{refgray}数据：JPM测算成本节约对核心EBIT净贡献约2亿英镑\par}
{\small\color{refgray}数据：市场共识已包含2026-2030年间约40个基点利润率提升\par}
{\small\color{refgray}边际变化：此前市场关注GSK成本节约计划本身，但JPM指出市场已部分预期利润率改善，额外增量有限，焦点转向管线数据兑现。\par}
\subsubsection*{关键数据}
\begin{itemize}
  \item 诺华中国区销售同比增长12\% vs 罗氏药品销售下滑33\%
  \item 17家未上市寿险公司核心偿付能力低于100\%，4家低于70\%
  \item 2025年55\%的财险公司出现承保亏损
  \item GSK成本节约对核心EBIT净贡献约2亿英镑，增量利润率约20个基点
  \item 保险深度从2015年3.5\%升至2025年4.4\%
  \item 前5大寿险公司市场份额45\%，前5大非寿险公司市场份额68\%
\end{itemize}
\begin{figure}[htbp]
\centering
\includegraphics[width=0.88\linewidth]{figures/fig_106.jpg}
\caption{Figure 1 - Figure 1: China – time series data for licensed insurers Number of legal entities Number of legal entities}
\end{figure}
\begin{figure}[htbp]
\centering
\includegraphics[width=0.88\linewidth]{figures/fig_107.jpg}
\caption{Figure 2 - Figure 2: China Insurance – life insurers and P\&C insurers' premium trend Rmb in billions, \%}
\end{figure}
\begin{figure}[htbp]
\centering
\includegraphics[width=0.88\linewidth]{figures/fig_108.jpg}
\caption{Figure 3 - Figure 3: China life insurers – core solvency ratio distribution \% Figure 4: China non-life insurers – core solvency ratio distribution \#\# Risk management score and insurance risk rating}
\end{figure}
\begin{figure}[htbp]
\centering
\includegraphics[width=0.88\linewidth]{figures/fig_109.jpg}
\caption{Figure 3 - Figure 3: China life insurers – core solvency ratio distribution \% Figure 4: China non-life insurers – core solvency ratio distribution \#\# Risk management score and insurance risk rating We have illustrated the divergence in c}
\end{figure}
{\small\color{refgray}本节综合 3 篇研究；涉及 NOM、JPM。}
\newpage
\subsection{金融与资金流：科技去杠杆接近尾声，中间价模型指向政策缓冲}
\textbf{全球科技与半导体板块的去杠杆进程快于预期，进一步压缩空间有限，市场可能接近阶段性底部；同时，人民币中间价模型显示逆周期因子仍在发挥缓冲作用，政策意图与市场定价之间的张力区间值得关注。}
\subsubsection*{主流共识}
\begin{itemize}
  \item 科技与半导体板块的去杠杆已基本完成，进一步下行空间有限。
  \item 逆周期因子仍在中间价形成中发挥缓冲作用，政策意图与市场定价存在差距。
\end{itemize}
\subsubsection*{分歧与条件差异}
\begin{itemize}
  \item 去杠杆完成后市场支撑来源：部分观点认为存储价格和云资本开支是主要支撑，另一部分则关注保证金账户杠杆数据尚未确认。
  \item 中间价方向：模型预测指向6.77附近，但最终定价取决于逆周期调节力度和外部环境变化。
\end{itemize}
\subsubsection*{机构观点}
\paragraph{JPM} 全球科技与半导体板块的去杠杆进程快于预期，目前进一步去杠杆的空间已相当有限。对冲基金在半导体指数（SOX）上的杠杆头寸已基本回吐，CTA等动量交易者的多头仓位大幅削减，此前流行的做空中国互联网+做多韩台日的相对交易也已完全反转。杠杆ETF因波动率损耗，其AUM占市值比已从6月峰值回落85\%（半导体非存储）和55\%（存储），不再是主要波动源。散户看涨期权买量已从6月5日近1400万张的峰值回落至历史均值附近。风险平价基金杠杆指标已正常化。唯一尚未确认的是保证金账户杠杆，官方数据要到8月底才出。后续支撑来自存储价格仍在上涨以及美国超大规模云厂商2026/2027年资本开支预期被大幅上调。
{\small\color{refgray}数据：杠杆ETF AUM占市值比从6月峰值回落85\%（半导体非存储）和55\%（存储）\par}
{\small\color{refgray}数据：散户看涨期权买量从6月5日近1400万张峰值回落至历史均值附近\par}
{\small\color{refgray}数据：对冲基金在SOX上的杠杆头寸已基本回吐\par}
{\small\color{refgray}数据：CTA等动量交易者在纳指、韩、台、日经上的多仓大幅削减\par}
{\small\color{refgray}边际变化：去杠杆速度比此前预期更快，目前进一步空间有限，市场可能接近阶段性底部。\par}
\paragraph{NOM} NOM亚洲外汇策略团队更新了美元兑人民币中间价预测模型，模型预测值从6.7899下调至6.7536，较前次官方收盘价低约119个基点。引入逆周期因子后的预测值为6.7727，较前次中间价低约172个基点，逆周期因子对模型预测的修正幅度约为191个基点，表明政策层面仍在通过该工具对中间价形成施加缓冲。模型误差近期呈现系统性偏大特征，显示模型输入变量与实际市场定价之间存在结构性错位。报告还列出了几个关键时间节点：7月底经济工作会议、8月中央行二季度货币政策报告、9月24日高层访美、9月底央行货币政策委员会会议。
{\small\color{refgray}数据：模型预测值6.7536，较前次下调363个基点\par}
{\small\color{refgray}数据：逆周期因子调整后预测值6.7727，较前次中间价低172个基点\par}
{\small\color{refgray}数据：逆周期因子修正幅度约191个基点\par}
{\small\color{refgray}数据：模型预测较前次官方收盘价低约119个基点\par}
{\small\color{refgray}边际变化：模型预测值较前次显著下调，逆周期因子修正幅度扩大，政策缓冲力度增强。\par}
\subsubsection*{关键数据}
\begin{itemize}
  \item 杠杆ETF AUM占市值比从6月峰值回落85\%（半导体非存储）和55\%（存储）
  \item 散户看涨期权买量从6月5日近1400万张峰值回落至历史均值附近
  \item 逆周期因子对中间价模型修正约191个基点
  \item 模型预测值6.7536，较前次下调363个基点
  \item 模型预测较前次官方收盘价低约119个基点
\end{itemize}
\begin{figure}[htbp]
\centering
\includegraphics[width=0.88\linewidth]{figures/fig_121.jpg}
\caption{Figure 1 - Figure 2: Short interest on Semiconductors ex Memory \& Hyperscalers individual stocks Short Interest as a \% share of shares outstanding.}
\end{figure}
\begin{figure}[htbp]
\centering
\includegraphics[width=0.88\linewidth]{figures/fig_123.jpg}
\caption{Figure 2 - Figure 4: Momentum signals for major equity indices}
\end{figure}
\begin{figure}[htbp]
\centering
\includegraphics[width=0.88\linewidth]{figures/fig_122.jpg}
\caption{Figure 1 - Figure 2: Short interest on Semiconductors ex Memory \& Hyperscalers individual stocks Short Interest as a \% share of shares outstanding. Figure 3: Short interest on SMH and DRAM US Equity ETFs Short Interest as a \% share of shar}
\end{figure}
\begin{figure}[htbp]
\centering
\includegraphics[width=0.88\linewidth]{figures/fig_124.jpg}
\caption{Figure 4 - Figure 4: Momentum signals for major equity indices Figure 5: Momentum signals for Asian equity indices weeks. As a result, the reduction in the AUM of leveraged equity ETFs has taken place much faster than we previously envisag}
\end{figure}
{\small\color{refgray}本节综合 2 篇研究；涉及 NOM、JPM。}
\newpage
\subsection{北美四轮车市场盈利修复与关税退税的参照效应}
\textbf{Polaris二季度财报显示，北美四轮车市场定价环境与成本结构正在改善，关税退税并非利润改善的唯一驱动。这一信号对川崎重工等同类企业具有参照意义，若退税效应落地，其PS\&E板块利润存在向上修正空间。}
\subsubsection*{主流共识}
\begin{itemize}
  \item 北美ORV（多功能越野车）零售市场整体呈低个位数增长，但头部企业通过产品组合优化和净定价改善实现利润修复。
  \item 关税退税（IEEPA）对部分企业利润形成一次性正面贡献，但剔除后仍可看到盈利基础改善。
  \item 促销费用控制是下半年利润能否超预期的关键变量，新车型推广顺利有望降低费用率。
\end{itemize}
\subsubsection*{分歧与条件差异}
\begin{itemize}
  \item Polaris已将关税退税纳入全年指引，而川崎重工尚未计入，若退税落地则存在额外增量。
  \item Polaris连续五个季度实现ORV市场份额提升，但川崎重工在摩托车和Side-by-side领域的市占率走势尚需观察。
  \item 市场对下半年北美四轮车销量增速的判断存在分歧，Polaris预计同比持平，但促销费用可能下降。
\end{itemize}
\subsubsection*{机构观点}
\paragraph{GS} Polaris二季度调整后销售额20.23亿美元，同比+9\%，有机增长+17\%。即便剔除7400万美元IEEPA关税退税，调整后EBITDA利润率仍同比回升540个基点至11.8\%，利润与利润率均实现同比增长。利润改善来自销量提升、产品组合优化、净定价改善及退税四个因素，其中退税并非唯一驱动。Polaris已连续五个季度实现北美ORV市场份额提升，零售单位销量同比+5\%，而整体市场仅低个位数增长。GS认为，这对川崎重工PS\&E板块具有参照意义：目前川崎重工1Q3/27利润预估70亿日元未计入任何退税影响，若退税落地则存在向上修正空间；同时，若新车型推广顺利，促销费用可能低于预期，全年PS\&E利润（研报估260亿日元，公司指引300亿日元）有上调可能。
{\small\color{refgray}数据：Polaris二季度调整后销售额20.23亿美元，同比+9\%，有机增长+17\%\par}
{\small\color{refgray}数据：剔除7400万美元关税退税后，调整后EBITDA利润率仍同比回升540个基点至11.8\%\par}
{\small\color{refgray}数据：Polaris ORV零售单位销量同比+5\%，整体市场仅低个位数增长\par}
{\small\color{refgray}数据：川崎重工1Q3/27 PS\&E利润预估70亿日元，未计入任何IEEPA退税\par}
{\small\color{refgray}边际变化：此前市场认为Polaris利润改善主要依赖一次性关税退税，但最新财报显示剔除退税后利润与利润率仍同比改善，表明北美四轮车市场定价环境与成本结构正在实质性修复。同时，川崎重工尚未将退税纳入预估，若退税落地将形成额外增量。\par}
\subsubsection*{关键数据}
\begin{itemize}
  \item Polaris二季度调整后销售额20.23亿美元，同比+9\%，有机增长+17\%
  \item 剔除7400万美元关税退税后，调整后EBITDA利润率仍同比回升540个基点至11.8\%
  \item Polaris ORV零售单位销量同比+5\%，整体市场仅低个位数增长
  \item Polaris已连续五个季度实现北美ORV市场份额提升
  \item 川崎重工1Q3/27 PS\&E利润预估70亿日元，未计入任何IEEPA退税
  \item 川崎重工全年PS\&E利润研报估260亿日元，公司指引300亿日元
\end{itemize}
\begin{figure}[htbp]
\centering
\includegraphics[width=0.88\linewidth]{figures/fig_092.jpg}
\caption{Exhibit 1 - Exhibit 2: Polaris: Polaris Powersports \& Off-Road segment revenue and gross margin trends}
\end{figure}
\begin{figure}[htbp]
\centering
\includegraphics[width=0.88\linewidth]{figures/fig_095.jpg}
\caption{Exhibit 3 - Exhibit 6: KHI: Side-by-side yoy retail unit trends (moving annual total) vs. Polaris ORV yoy retail unit momentum}
\end{figure}
\begin{figure}[htbp]
\centering
\includegraphics[width=0.88\linewidth]{figures/fig_093.jpg}
\caption{Exhibit 1 - Exhibit 2: Polaris: Polaris Powersports \& Off-Road segment revenue and gross margin trends Exhibit 3: Polaris: Adjusted EBITDA bridge Exhibit 4: Polaris: Inventory levels and inventory turnover (months of sales)}
\end{figure}
\begin{figure}[htbp]
\centering
\includegraphics[width=0.88\linewidth]{figures/fig_094.jpg}
\caption{Exhibit 2 - Exhibit 2: Polaris: Polaris Powersports \& Off-Road segment revenue and gross margin trends Exhibit 3: Polaris: Adjusted EBITDA bridge Exhibit 4: Polaris: Inventory levels and inventory turnover (months of sales) Exhibit 5: KH}
\end{figure}
{\small\color{refgray}本节综合 1 篇研究；涉及 GS。}
\newpage
\subsection{科技与工业：设备需求强劲，消费与化工转型分化}
\textbf{半导体设备与存储需求在产能约束下持续超预期，但消费娱乐与化工行业面临结构性转型与现金流压力，行业间分化明显。}
\subsubsection*{主流共识}
\begin{itemize}
  \item 半导体设备需求强劲，Lam Research营收与指引均超预期，WFE预期上调至低1500亿美元区间。
  \item 企业级HDD供需紧张，Seagate近线HDD供应已分配至2028年，客户开始谈判2029年合同。
  \item 化工行业销量回升但价格偏弱，下游补库驱动而非终端需求实质性改善。
\end{itemize}
\subsubsection*{分歧与条件差异}
\begin{itemize}
  \item Lam Research毛利率创20年新高至52\%，长期目标上调至55\%，但中国收入占比从34\%降至26\%，区域风险仍在。
  \item Seagate出货量同比增速超40\%，但公司维持25\%的长期CAGR目标，扩产依赖技术迁移而非增加数量。
  \item 迪士尼主题乐园客流增长放缓，Citi认为Q4市场共识的6\%客流增长存在不确定性，管理层可能转向全球客流指标。
  \item 巴斯夫运营现金流大幅下滑，现金转化率从89\%降至21\%，转型支出与营运资本占用是主因，但全年指引仍维持。
\end{itemize}
\subsubsection*{机构观点}
\paragraph{Bernstein} Lam Research交出‘最好财报之一’，营收67.2亿美元、每股收益1.82美元均超预期，下一季度营收指引81亿美元，较市场共识高约10亿美元。毛利率52\%创20年新高，受益于定价、运营效率与产品组合，长期毛利率目标上调至55\%。NAND业务环比增长118\%，逻辑增长63\%，但中国收入占比从34\%降至26\%，预计趋势延续。WFE预期从1400亿上调至低1500亿美元，先进封装收入预计增长超70\%。
{\small\color{refgray}数据：营收67.2亿美元，超预期66.8亿\par}
{\small\color{refgray}数据：每股收益1.82美元，超预期1.70\par}
{\small\color{refgray}数据：毛利率52\%，创20年新高\par}
{\small\color{refgray}数据：NAND业务环比增长118\%\par}
{\small\color{refgray}边际变化：毛利率突破52\%并上调长期目标至55\%，显示定价能力与产品结构发生结构性变化，而非单纯季度波动。\par}
\paragraph{NOM} Seagate 26/6财年Q4营收36.29亿美元，环比增长17\%，营业利润率44.6\%，均超预期。Q1指引营收中值41亿美元，营业利润率约50\%，同样大幅高于市场预测。近线HDD的2028年供应已全部被客户锁定，客户已开始谈判2029年合同，反映AI与云服务商长期需求重塑采购模式。尽管出货量同比增速超40\%，公司仍维持25\%的长期CAGR目标，扩产依赖技术迁移。每Exabyte价格同比上涨超10\%，定价能力增强。
{\small\color{refgray}数据：Q4营收36.29亿美元，超预期34.9亿\par}
{\small\color{refgray}数据：营业利润率44.6\%，超预期41.9\%\par}
{\small\color{refgray}数据：Q1营收指引中值41亿美元\par}
{\small\color{refgray}数据：近线HDD供应已分配至2028年\par}
{\small\color{refgray}边际变化：客户采购周期从按需补货转向长期锁定产能，2028年供应已分配完毕，2029年合同谈判启动，传统12-18个月锁定周期被打破。\par}
\paragraph{Citi} 迪士尼2026财年调整后每股收益预测6.80美元，略低于公司16\%增长指引。Q3分部EBIT预计52.35亿美元，低于公司53亿指引，娱乐板块EBIT比市场共识低10.5\%，流媒体EBIT增速85.5\%低于预期90.9\%。主题乐园EBIT 29.54亿美元，同比增长17.4\%，成为主要支撑。但Q4市场共识的6\%客流增长存在不确定性，管理层可能将指标从国内客流转向全球客流。SVOD内容支出增加，下调FY27/FY28 EBIT利润率预期130bps和200bps。
{\small\color{refgray}数据：2026财年Adj EPS预测6.80美元\par}
{\small\color{refgray}数据：Q3分部EBIT预计52.35亿美元\par}
{\small\color{refgray}数据：娱乐板块EBIT比共识低10.5\%\par}
{\small\color{refgray}数据：流媒体EBIT增速85.5\%\par}
{\small\color{refgray}边际变化：市场焦点从单一增长叙事转向多业务线同步调整，主题乐园客流增长放缓与流媒体内容支出上升成为新关注点。\par}
\paragraph{GS} 巴斯夫Q2调整后EBITDA 24.49亿欧元，超预期18\%，销售额172.06亿欧元，超预期5\%。销量同比+7.3\%，远超预期+2.0\%，但价格仅+11.5\%，低于预期+12.7\%，呈现‘量强价弱’组合，反映下游补库驱动。利润结构分化，材料和工业解决方案EBITDA分别超预期28\%和29\%，但化学品板块利润率低486bps。运营现金流仅5.24亿欧元，远低于去年同期的15.85亿，现金转化率从89\%骤降至21\%，转型支出与营运资本增加是主因。
{\small\color{refgray}数据：调整后EBITDA 24.49亿欧元，超预期18\%\par}
{\small\color{refgray}数据：销售额172.06亿欧元，超预期5\%\par}
{\small\color{refgray}数据：销量同比+7.3\%，远超预期+2.0\%\par}
{\small\color{refgray}数据：集团价格仅+11.5\%，低于预期+12.7\%\par}
{\small\color{refgray}边际变化：运营现金流大幅下滑揭示转型真实成本，现金转化率骤降，但全年EBITDA指引69-77亿欧元维持不变。\par}
\subsubsection*{关键数据}
\begin{itemize}
  \item Lam Research营收67.2亿美元，超预期66.8亿
  \item 毛利率52\%创20年新高，长期目标上调至55\%
  \item Seagate近线HDD供应已分配至2028年
  \item 每Exabyte价格同比上涨超10\%
  \item 迪士尼Q3分部EBIT预计52.35亿美元，低于公司指引
  \item 流媒体EBIT增速85.5\%，低于预期90.9\%
  \item 巴斯夫销量同比+7.3\%，远超预期+2.0\%
  \item 运营现金流5.24亿，去年同期15.85亿
\end{itemize}
\begin{figure}[htbp]
\centering
\includegraphics[width=0.88\linewidth]{figures/fig_023.jpg}
\caption{EXHIBIT 2 - EXHIBIT 4: NAND and Logic saw strong growth QoQ, while Foundry and DRAM were down slightly QoQ EXHIBIT 5: Memory share was up 7ppt QoQ to 46\% with equal share from DRAM and NAND Lam Research}
\end{figure}
\begin{figure}[htbp]
\centering
\includegraphics[width=0.88\linewidth]{figures/fig_024.jpg}
\caption{EXHIBIT 3 - EXHIBIT 5: Memory share was up 7ppt QoQ to 46\% with equal share from DRAM and NAND Lam Research Memory Contribution to Systems Revenues}
\end{figure}
\begin{figure}[htbp]
\centering
\includegraphics[width=0.88\linewidth]{figures/fig_074.jpg}
\caption{Figure 5 - \#\# \#2: US Attendance Trends Over the past few quarters, investors have closely followed Domestic attendance. To gauge recent trends, we looked at three external data sources: ■ First, in F3Q26, Orlando airport passengers (through May) grew 3\% YoY. Figure 5. Or}
\end{figure}
\begin{figure}[htbp]
\centering
\includegraphics[width=0.88\linewidth]{figures/fig_076.jpg}
\caption{Figure 6 - © 2026 Citi Inc. No redistribution without Citi's written permission. ■ Second, Placer foot traffic at Disney’s Domestic parks is up 5\% YoY in F3Q26. Figure 6. US Parks Avg Daily Foot Traffic (000s; percent) © 2026 Citi Inc. No redistribution without Citi's wr}
\end{figure}
{\small\color{refgray}本节综合 4 篇研究；涉及 Bernstein、NOM、Citi、GS。}
\newpage
\subsection{中国经济、地产与金融：AI驱动的区域K型分化加剧}
\textbf{中国经济增长正经历结构性转变，AI产业红利高度集中于少数智慧城市，导致地理层面的K型分化加剧，而占GDP大部分的其他地区增速放缓，这一格局削弱了全面刺激政策的必要性。}
\subsubsection*{主流共识}
\begin{itemize}
  \item 中国经济增长动能出现结构性分化，AI相关活动是主要驱动力。
  \item 智慧城市（一线及精选二线）的经济表现显著优于全国其他地区。
  \item 全国固定资产投资整体收缩，但智慧城市（如北京、上海）逆势增长。
  \item 总量GDP增速放缓掩盖了区域间的根本性重组。
\end{itemize}
\subsubsection*{分歧与条件差异}
\begin{itemize}
  \item 关于AI增长红利的扩散范围：部分观点认为AI红利将逐步扩散至更多城市，而NOM报告指出目前仅集中在少数智慧城市。
  \item 政策应对方向：若K型分化持续，全面刺激的必要性降低，但结构性支持政策（如针对智慧城市的产业政策）可能加码。
\end{itemize}
\subsubsection*{机构观点}
\paragraph{NOM} NOM指出，中国经济增长正出现地理层面的K型分化，AI驱动的增长红利主要被少数智慧城市捕获。2026年上半年，七个智慧城市（三个一线+四个精选二线）的加权平均实际GDP增速从2025年的5.4\%升至5.6\%，而占全国GDP约83\%的其他地区增速从4.9\%降至4.5\%。智慧城市贡献了全国上半年增长的20\%，为至少二十年来最高。固定资产投资数据进一步佐证：全国FAI上半年同比收缩5.7\%，但北京和上海分别实现3.0\%和6.8\%的正增长。报告认为，AI产业集聚效应正在重塑城市间增长格局，而非均匀扩散，且少数城市的强劲增长难以抵消大部分地区的放缓，全面刺激政策的可能性降低。
{\small\color{refgray}数据：全国实际GDP增速从2025年的5.0\%降至2026年上半年的4.7\%。\par}
{\small\color{refgray}数据：七个智慧城市加权平均增速从5.4\%升至5.6\%。\par}
{\small\color{refgray}数据：其他地区（占GDP约83\%）增速从4.9\%降至4.5\%。\par}
{\small\color{refgray}数据：智慧城市贡献全国上半年增长的20\%，为至少二十年最高。\par}
{\small\color{refgray}边际变化：此前市场普遍关注总量GDP增速放缓，但NOM首次系统量化了AI驱动的区域K型分化，并指出智慧城市增长加速而其他地区降幅更大，这一结构性变化改变了政策预期（全面刺激必要性降低）。\par}
\subsubsection*{关键数据}
\begin{itemize}
  \item 全国实际GDP增速从2025年的5.0\%降至2026年上半年的4.7\%。
  \item 七个智慧城市加权平均增速从5.4\%升至5.6\%。
  \item 其他地区（占GDP约83\%）增速从4.9\%降至4.5\%。
  \item 智慧城市贡献全国上半年增长的20\%，为至少二十年最高。
  \item 全国FAI上半年同比收缩5.7\%。
  \item 北京FAI同比增长3.0\%，上海FAI同比增长6.8\%。
\end{itemize}
{\small\color{refgray}本节综合 1 篇研究；涉及 NOM。}
\newpage
\section{辅助信号｜战略咨询}
本板块整合波士顿咨询近期关于AI现场服务、生物制药剥离、资产剥离退出路径、替代性交易及2025年并购趋势的研究，提炼出三个核心主题：AI驱动的现场服务转型、并购与剥离的执行分化、以及替代性交易的结构性崛起。
\subsection{AI驱动的现场服务转型：从隐性知识到系统化增长}
\textbf{现场服务行业面临技师短缺与知识断层的双重约束，AI智能体与生成式AI助手可将隐性经验转化为可复用系统，实现15\%以上收入增长和5个百分点以上毛利率提升。}
\begin{itemize}
  \item 美国卡车运输业因技师短缺年损24亿美元，工业制造商非计划停机损失达500亿美元/年；暖通空调需求增长15\%但供给仅增6\%，航空维修2025年技师缺口10\%。
  \item 技师平均年龄40-54岁，退休带走数十年“部落知识”，新技师需快速再培训；AI工具通过知识库化隐性经验，降低知识传递效率瓶颈。
  \item 设备联网、生成式AI助手与AI智能体三重技术变化同步发生，报告基于1000+客户合作提出从传统服务向AI增强型转型路径。
\end{itemize}
\begin{figure}[htbp]
\centering
\includegraphics[width=0.88\linewidth]{figures/fig_160.jpg}
\caption{Exhibit 1 - Exhibit 1 - Carve-Outs Offer Significant Value, but Performance Varies Widely Sources: S\&P Capital IQ; Refinitiv Datastream; BCG analysis. Note: The Bloomberg Spin-Off Index tracks US spinoffs for the first three years of their}
\end{figure}
\subsection{并购与剥离的执行分化：价值创造差距超100个百分点}
\textbf{资产剥离占并购市场比重持续上升至近50\%，但只有约50\%的交易中期创造价值；执行精细度决定成败，生物制药剥离面临五重独特复杂性。}
\begin{itemize}
  \item 2013年资产剥离占并购市场近50\%，2014上半年交易额同比跃升62\%，超350亿美元大型交易占35\%以上；剥离后EBITDA利润率平均提升超1个百分点。
  \item 生物制药剥离需应对全球布局与本地运营矛盾、监管环境动态变化、高度整合供应链、经验证IT系统分离及产品供应连续性等五重挑战，每笔交易可能触发数千项监管变更。
  \item 2020-2021年全球季度平均剥离交易价值3090亿美元，比五年均值高25\%；表现最佳与最差交易股东总回报差距超100个百分点，差异主要来自执行层面。
\end{itemize}
\begin{figure}[htbp]
\centering
\includegraphics[width=0.88\linewidth]{figures/fig_163.jpg}
\caption{Exhibit 1 - After furrowing a deep trough in the years following 2008, the M\&A market has at last recovered to the levels of 2005 and 2006. That said, most market participants saw 2013 as a disappointment. Overall deal volume declined by 6 percent and total deal value fel}
\end{figure}
\begin{figure}[htbp]
\centering
\includegraphics[width=0.88\linewidth]{figures/fig_161.jpg}
\caption{Exhibit 1 - Exhibit 2 - Biopharma Divestitures Reached an All-Time High in 2019}
\end{figure}
\subsection{替代性交易的结构性崛起：少数股权与联盟成为常规工具}
\textbf{少数股权、合资与战略联盟等替代性交易占比从1990年的20\%-25\%升至约35\%，在扰动期占比加速攀升，成为企业获取技术能力与应对不确定性的常规选择。}
\begin{itemize}
  \item 2019年全球合资与联盟交易量达1.1万笔历史新高，软件与IT服务、医疗设备领域增长最显著；2020年上半年整体交易量骤降80\%，但少数股权交易占比反而走高。
  \item 驱动因素中54\%受访者指向技术变革，54\%指向商业模式变化，45\%提及不确定性分担；替代性交易公告回报率呈上升趋势，但中长期表现不到一半交易创造价值。
  \item 2025年前九个月全球交易总额同比增10\%，但亚太跌至十年低点；北美占62\%全球交易价值，27宗超百亿美元超级交易高于2024年同期21宗，但大型交易数量仍处历史低位。
\end{itemize}
\begin{figure}[htbp]
\centering
\includegraphics[width=0.88\linewidth]{figures/fig_168.jpg}
\caption{EXHIBIT 1 - tions helped to drive an 11\% increase in deal value in the region compared with 2018, while deal volume decreased by 10\%. In contrast, only five megadeals were announced in Europe, two of which were withdrawn. The scarcity of megadeals in Europe contributed to}
\end{figure}
\begin{figure}[htbp]
\centering
\includegraphics[width=0.88\linewidth]{figures/fig_173.jpg}
\caption{EXHIBIT 1 - Private equity (PE) and venture capital (VC) activity is trending higher. Global PE deal value rose 38\% year-to-date through the first three quarters of 2025 compared with the same period last year, driven by large deal activity such as the mega buyout of Elec}
\end{figure}
\newpage
\section{辅助信号｜智库/国际机构}
本日新增报告涵盖亚洲人力发展不平等、新兴市场房地产数据基建、AI对竞争格局的影响、贸易争端机制、教育干预效果及公民参与政策设计等议题，为宏观与产业叙事提供实证信号。
\subsection{人力资本与教育干预的实证路径}
\textbf{亚洲人力发展不平等损失高达35\%，教育是最大短板；加纳阅读干预项目九个月提升成绩0.5个标准差，证明结构化教学与差异化指导的可扩展性。}
\begin{itemize}
  \item 亚洲开发银行报告显示，亚洲因不平等导致的人力发展指数损失在9.7\%至35.3\%之间，阿富汗损失最高（35.3\%），文莱最低（9.7\%）；教育不平等是主要驱动因素，阿富汗（48.8\%）、不丹（48.2\%）教育不平等程度最高。
  \item 性别不平等与人力发展损失呈正相关，阿富汗、巴基斯坦、孟加拉国位于高脆弱性象限；新加坡、中国、马来西亚位于强表现者象限，性别平等指数与政府效能系统性关联。
  \item 经合组织AI辅助细粒度分析显示，泰国15岁学生PISA 2022数学得分393.9分（低于OECD平均472.4分），优势在几何、测量、估算等程序性任务，短板在数学化情境、结果解释与函数关系识别。
  \item 世界银行加纳阅读干预项目（TFLI）九个月使一年级学生阅读成绩提升0.504个标准差，相当于超过两年学习进步；效果排全球同类项目前9\%，基础技能（字母发音）效应量达0.764个标准差。
  \item TFLI项目采用智能手机增强的结构化教学法，通过半脚本化教案、学生练习册和SmartCoach应用实现差异化教学，模型预测基础技能制约高阶技能，实证结果与模型高度一致。
\end{itemize}
\begin{figure}[htbp]
\centering
\includegraphics[width=0.88\linewidth]{figures/fig_139.jpg}
\caption{Figure 2 - Figure 1: Human Development Inequality Across Asian Economies A. HDI Loss due to Inequality (\%), 2023 B. Human Inequality Coefficient, 2023 HDI = human development index, Lao PDR = Lao People's Democratic Republic; PRC = Peopl}
\end{figure}
\begin{figure}[htbp]
\centering
\includegraphics[width=0.88\linewidth]{figures/fig_141.jpg}
\caption{Figure 2 - Figure 2: Dimensions of Human Development Inequality Across Asian Economies A. Inequality in Life Expectancy, 2023 B. Inequality in Education, 2023 C. Inequality in Income, 2023 PRC = People's Republic of China. Note: Panel}
\end{figure}
\subsection{AI对竞争格局的动态影响}
\textbf{非生成式AI未显著推高企业市场势力，但生成式AI可能同时为小企业创造机会并强化已有优势企业地位，AI专利活动与利润率增长正相关。}
\begin{itemize}
  \item 经合组织基于法国、葡萄牙2011-2022年企业数据发现，非生成式AI采用者规模更大、生产率更高，但利润率未显著高于非采用者，未系统性推动企业向上移动市场份额分布。
  \item 生成式AI扩散可能加剧分化：AI创新领域集中度与销售集中度正相关，AI专利活动与更高利润率增长相关，尤其在ICT行业。
  \item 报告构建了区分AI使用者与开发者、生成式与非生成式AI的分析框架，指出当前AI采用率仍较低（约20\%），竞争效应尚未完全显现，需持续监测。
\end{itemize}
\begin{figure}[htbp]
\centering
\includegraphics[width=0.88\linewidth]{figures/fig_150.jpg}
\caption{Figure 4 - The key independent variable is \$AI\_\{i,t\}\$ , a dummy variable that indicates whether a firm uses AI or not. \$\textasciicircum{}\{36\textbackslash{} 37\}\$ \$X\_\{it\}\$ accounts for the use of fast broadband (broadband connection with at least 100 Mbit/s) to capture digital infrastructure and the sh}
\end{figure}
\begin{figure}[htbp]
\centering
\includegraphics[width=0.88\linewidth]{figures/fig_153.jpg}
\caption{Figure 4 - Panel A of Figure 4.3 reports the trends of these measures over the years 2001-21. The early years of the sample display relatively high concentration, likely reflecting the presence of only a limited number of AI innovators. The various indicators decline unt}
\end{figure}
\subsection{贸易争端与数据基建的制度信号}
\textbf{巴西就美国301关税诉诸WTO争端解决机制，触及多边规则核心；IMF技术援助推动科索沃与乌克兰住宅价格指数编制，强化新兴市场数据基础设施。}
\begin{itemize}
  \item 巴西正式向WTO起诉美国依据301条款加征的25\%和12.5\%额外关税，主张违反GATT最惠国待遇和DSU第23条禁止单边报复规则；两项关税叠加后巴西产品进入美国需额外缴纳37.5\%从价关税。
  \item 巴西诉讼切入点精准：不挑战301调查事实认定，而是攻击程序合法性——美国未诉诸WTO争端解决机制而自行采取关税行动，触及多边贸易规则核心。
  \item IMF技术援助显示，科索沃首个官方住宅价格指数计划2025年底前发布，基于地籍局行政数据（2018年后记录可靠），采用特征价格模型加四季度滚动窗口。
  \item 乌克兰住宅价格指数基于挂牌数据编制，回归模型R平方值大多超过0.5，已覆盖八个季度公寓数据；行政数据因申报价格显著低于要价，暂无法用于官方编制。
\end{itemize}
\subsection{AI、技术与生产率}
\textbf{用于判断 AI 投资周期、生产率扩散和产业约束是否正在改变资产定价。}
\begin{itemize}
  \item 欧盟凝聚政策是其最大的研究战略，每年投入数百亿欧元用于缩小区域发展差距。但一个长期存在的问题是：这些资金是否真正反映了当地居民的需求？经合组织与欧盟委员会区域与城市政策总司合作，在7个欧盟国家的11个地区设计了不同类型的公民参与试点，试图回答这个问题。这份报告的核心判断是：当参与...
\end{itemize}
\begin{figure}[htbp]
\centering
\includegraphics[width=0.88\linewidth]{figures/fig_149.jpg}
\caption{Figure 1 - \textbackslash{}- Contribute to the ongoing reform of Cohesion Policy by showing the benefits of the Partnership Principle and extending the involvement of citizens and civil society as partners of Cohesion Policy. A first phase of this collaboration implemented in 2020 and }
\end{figure}
\section{结语}
后续需验证的关键节点包括：Meta 2027-2028年资本支出路径的具体框架、存储定价上行空间的持续性（尤其HBM）、中国消费K型分化是否向更多品类扩散、美国电力变压器缺口收窄速度、以及AI对SaaS收入贡献能否从低个位数向拐点演进。
\newpage
\section{覆盖概览}
投行/券商 63 篇；战略咨询 5 篇；智库/国际机构 8 篇。\par
投行覆盖：GS 20篇 / Bernstein 13篇 / MS 11篇 / NOM 8篇 / JPM 4篇 / BARC 2篇 / HSBC 2篇 / UBS 2篇 / Citi 1篇\par
\section{Disclaimer}
{\scriptsize\color{refgray}
This community is a paid learning and information-sharing community created for general educational purposes only. The membership fee is charged solely for access to community discussions, educational content, organization, and operational costs. It is not an advisory fee, management fee, performance fee, brokerage fee, or compensation for personalized financial, investment, legal, tax, or professional advice.\par
I participate and share information solely as an individual and for educational discussion among community members. I am not acting as a financial adviser, investment adviser, broker, dealer, portfolio manager, legal adviser, tax adviser, accountant, fiduciary, or any other licensed professional adviser. Nothing shared in this community constitutes, and should not be interpreted as, financial advice, investment advice, legal advice, tax advice, a recommendation, solicitation, offer, endorsement, instruction, or guarantee to buy, sell, hold, short, trade, or invest in any security, cryptocurrency, fund, derivative, strategy, or other financial product.\par
All content shared in this community, including messages, discussions, links, files, charts, examples, market commentary, watchlists, AI-generated or AI-polished summaries, and third-party materials, is general, impersonal, and educational in nature. It does not take into account any individual's financial situation, investment objectives, risk tolerance, investment horizon, liquidity needs, tax status, legal restrictions, or suitability requirements. No content should be relied upon as suitable for any specific person, account, portfolio, or financial decision.\par
Information shared in this community may be incomplete, inaccurate, outdated, speculative, AI-generated, AI-edited, or based on third-party internet sources that have not been independently verified. I make no representation or warranty as to the accuracy, completeness, timeliness, reliability, or suitability of any information shared.\par
Financial markets involve substantial risk, including the possible loss of principal. Past performance, historical data, hypothetical examples, simulated results, backtests, personal opinions, or educational examples do not guarantee future results. Each member is solely responsible for conducting their own research, exercising independent judgment, and making their own decisions. Before making any financial, investment, legal, or tax decision, members should consult qualified and licensed professionals.\par
Participation in this community, payment of any membership fee, or communication with me or other members does not create any adviser-client, fiduciary, professional, contractual, agency, partnership, employment, or other advisory relationship. By joining or participating in this community, you acknowledge and agree that you are solely responsible for your own decisions, actions, transactions, profits, losses, risks, and consequences, and that I shall not be liable for any direct, indirect, incidental, consequential, financial, legal, tax, or other loss or damage arising from or related to any content, discussion, information, or material shared in this community.\par
}
\newpage
\begin{center}
{\Large 更多详情报告kcdesk.com}\par\vspace{0.8cm}
\includegraphics[width=0.58\linewidth]{\detokenize{../../prompts/zsxq_img.jpg}}
\end{center}
\end{document}
